\documentclass[11pt]{article}

\usepackage[a4paper,margin=1in]{geometry}
\usepackage[T1]{fontenc}
\usepackage{lmodern}
\usepackage{graphicx}
\usepackage{amsmath,amssymb}
\usepackage{booktabs}
\usepackage{array}
\usepackage{tabularx}
\usepackage{longtable}
\usepackage{float}
\usepackage{subcaption}
\usepackage{hyperref}
\usepackage{enumitem}
\usepackage{xcolor}
\usepackage{microtype}
\usepackage{parskip}

\hypersetup{
    colorlinks=true,
    linkcolor=blue,
    urlcolor=blue,
    citecolor=blue
}

\setlength{\parskip}{6pt}

\title{%
  \textbf{Wafer Map Anomaly Detection on WM-811K}\\[0.4em]
  \large From Reconstruction Baselines to Pretrained Local Feature Modelling
}
\author{Deep Learning Project Report, Group 08\\[0.3em]
  Henry Lee Jun (1004219) \quad Chia Tang (1007200) \quad Genson Low (1005931)}
\date{\today}

\begin{document}

\begin{titlepage}
  \centering
  \vspace*{2cm}
  {\LARGE\bfseries Wafer Map Anomaly Detection on WM-811K\par}
  \vspace{0.5cm}
  {\large From Reconstruction Baselines to Pretrained Local Feature Modelling\par}
  \vspace{1.5cm}
  {\large\bfseries Deep Learning Project Report\par}
  \vspace{0.4cm}
  {\large 50.039 Deep Learning, Y2026\par}
  \vspace{0.4cm}
  {\large Group 08\par}
  \vspace{0.6cm}
  {\normalsize
    Henry Lee Jun (1004219) \quad
    Chia Tang (1007200) \quad
    Genson Low (1005931)\par}
  \vspace{0.8cm}

  {\large \today\par}
  \vfill
  {\small All models implemented in PyTorch. This report documents the experimental design,
  results, and diagnostic analysis for normal-only wafer anomaly detection on the
  WM-811K / LSWMD dataset.\par}
\end{titlepage}

\pagenumbering{roman}

\begin{abstract}
This report presents an anomaly-detection study on the WM-811K / LSWMD wafer-map dataset
under a strict normal-only training regime. The project begins with convolutional autoencoders
to establish a stable baseline pipeline, then moves progressively toward pretrained
ImageNet representations, teacher-student distillation, and PatchCore-style local anomaly
scoring. The main finding is that stronger frozen backbones combined with local feature
modelling improve performance substantially over plain reconstruction or global
embedding-distance baselines. The best completed run on the main benchmark is a
direct-$224\!\times\!224$ PatchCore model built on a pretrained ViT-B/16 backbone
(block 6 features, top-10\% patch reduction), achieving $\mathrm{F1}=0.595$,
$\mathrm{AUROC}=0.956$, and $\mathrm{AUPRC}=0.671$ under a deployment-style
95th-percentile validation threshold.
Score histograms and UMAP diagnostics show that the remaining bottleneck is
score-distribution overlap and naive thresholding rather than feature quality alone.
\end{abstract}

\newpage
\tableofcontents
\newpage

\pagenumbering{arabic}

\section{Introduction and Motivation}
\label{sec:problem}

After fabrication, each die on a semiconductor wafer is electrically tested and
assigned a pass or fail label at its grid location, producing a wafer bin
map (WBM). These spatial maps are diagnostically valuable because different
failure patterns indicate different root causes. A ring of failures near the wafer
perimeter suggests a different process error than a central cluster or a diagonal
scratch. Automatically detecting such anomalies enables faster wafer disposition
and improves upstream process control.

Supervised multi-class classification is a natural first approach, and prior work
on WM-811K demonstrates that it can be highly effective in closed benchmarks.
Recent CNN and ensemble classifiers report classification accuracies exceeding
99\% on balanced splits drawn from the same eight labelled defect categories
\cite{khatun2026cbam, taha2025survey}. However, these numbers measure a
different task. The classifier is asked: given a wafer known to contain one of these eight defect
types, which one is it? Production anomaly detection asks instead:
given a randomly drawn wafer from the line, is it defective at all?
These are fundamentally different problems. The high benchmark accuracy arises
because the classifier operates only on confirmed defective wafers drawn from
known categories. In production the vast majority of wafers are normal, and
entirely new failure modes can appear without warning.

Beyond this task mismatch, supervised classification faces three compounding
practical limitations in real fabrication environments:

\begin{itemize}[leftmargin=1.5em]
  \item \textbf{Closed-world constraint.} A classifier with 99\% accuracy on
        eight known defect types has exactly 0\% recall on any new failure mode
        introduced by a process change or new equipment. The defect vocabulary
        in a real fab is neither fixed nor complete.
  \item \textbf{Severe label scarcity.} Only $3.1\%$ of WM-811K wafers carry a
        confirmed defect label and $78.7\%$ are entirely unlabelled. Labelling
        a defect requires expert engineers to correlate wafer maps with equipment
        logs, typically days or weeks after the wafer leaves the line.
  \item \textbf{Scaling cost.} Adding a new defect class requires repeating the
        full label-acquisition cycle from scratch. A normal-only system needs
        only verified-clean wafers, which are abundant and cheap to collect.
\end{itemize}

This project therefore adopts normal-only anomaly detection. The model
is trained exclusively on the 147,431 labelled-normal wafers. Any wafer that
deviates from the learned normal representation receives a high anomaly score,
regardless of which defect type caused the deviation. No defect labels are
required at any stage: training, threshold selection, or deployment. The
deployment threshold is the 95th percentile of validation-normal scores,
matching the production constraint that only normal wafers are available for
calibration. Because false negatives (missed defects) are the more critical
error in semiconductor inspection, this report prioritises F1 at the deployed
threshold rather than AUROC alone.

\section{Dataset and Evaluation Protocol}

\subsection{Dataset}

The WM-811K dataset contains 811,457 wafer bin maps (Table~\ref{tab:dataset_distribution}).
Each map is converted to a normalised float array in $[0,1]$ and resized using
nearest-neighbour interpolation. Only explicitly labelled-normal wafers are used
for training. All labelled defect types are treated as anomaly at evaluation time.
The unlabelled $78.7\%$ is excluded entirely to prevent label contamination.

\begin{table}[H]
\centering
\caption{WM-811K dataset composition.}
\label{tab:dataset_distribution}
\begin{tabular}{lcc}
\toprule
Category & Count & Fraction \\
\midrule
No-label & ${\approx}638{,}507$ & $78.7\%$ \\
Labelled: None (Normal) & $147{,}431$ & $18.2\%$ \\
Labelled: Pattern (Defect) & $25{,}519$ & $3.1\%$ \\
\midrule
\textbf{Total} & \textbf{811,457} & \textbf{100\%} \\
\bottomrule
\end{tabular}
\end{table}

Most experiments resize wafer maps to $64\!\times\!64$ pixels. Later
pretrained-backbone experiments additionally evaluate direct $224\!\times\!224$
preprocessing to match each backbone's native resolution.

\subsection{Benchmark Split}

All experiments share a single split drawing 50,000 labelled-normal wafers,
with defects held out exclusively in the test partition (Table~\ref{tab:main_split}).
Normals are divided 80/10/10 across train, validation, and test.
Defects are capped at 5\% of test-normal count to reflect realistic production
prevalence. The validation set serves one purpose only: deriving the deployment
threshold. No defect labels are visible before test evaluation.

\begin{table}[H]
\centering
\caption{Main benchmark split.}
\label{tab:main_split}
\begin{tabular}{lcc}
\toprule
Subset & Normal wafers & Defect wafers \\
\midrule
Train   & 40{,}000 & 0   \\
Val     & 5{,}000  & 0   \\
Test    & 5{,}000  & 250 \\
\bottomrule
\end{tabular}
\end{table}

\subsection{Defect Categories}

The 250 test-set defects span eight pattern types that vary substantially in
spatial scale and coverage (Table~\ref{tab:defect_types}, Figure~\ref{fig:defect_gallery}).

\begin{table}[H]
\centering
\caption{Defect-type breakdown of the main benchmark test set.}
\label{tab:defect_types}
\begin{tabularx}{\linewidth}{lc>{\raggedright\arraybackslash}X}
\toprule
Defect type & Count & Spatial character \\
\midrule
Edge-Ring & 84  & Continuous ring of failures near the wafer perimeter \\
Edge-Loc  & 53  & Localised cluster at or near the wafer edge \\
Center    & 50  & Compact cluster at the wafer centre \\
Loc       & 34  & Off-centre localised patch \\
Scratch   & 15  & Thin elongated streak across the die grid \\
Donut     & 7   & Curved partial ring inside the wafer body \\
Random    & 5   & Scattered failures with no spatial coherence \\
Near-full & 2   & Defective region covering nearly the entire wafer \\
\bottomrule
\end{tabularx}
\end{table}

\begin{figure}[H]
    \centering
    \includegraphics[width=\linewidth]{images/defect_type_gallery.png}
    \caption{Representative examples of all eight defect pattern types in the
    WM-811K evaluation set. Each image is a wafer bin map rendered as a
    binary die-level pass/fail grid. The diversity in spatial scale and
    structure, from full-wafer ring patterns to single-streak scratches,
    is a primary challenge for any global anomaly scoring approach. Generated by \texttt{experiments/anomaly\_detection/report\_figures/notebook.ipynb}.}
    \label{fig:defect_gallery}
\end{figure}

Broad defects such as Edge-Ring and Center are detectable by any model that
measures overall reconstruction quality. Small, localised defects such as
Scratch and Loc require spatially precise, patch-level scoring. This
distinction drives the experimental progression.

\section{Methods}
\label{sec:methodology}

Table~\ref{tab:method_overview} summarises all model families. The progression
moves from fully in-domain reconstruction models through frozen pretrained
backbones to local patch-level scoring, testing four questions along the way.

\begin{table}[H]
\centering
\caption{Overview of model families evaluated in this study.}
\label{tab:method_overview}
\begin{tabularx}{\linewidth}{>{\raggedright\arraybackslash}p{3.2cm}
                              >{\raggedright\arraybackslash}p{2.2cm}
                              >{\raggedright\arraybackslash}p{2.4cm}
                              >{\raggedright\arraybackslash}X}
\toprule
Method & Training & Scoring & Backbones / notes \\
\midrule
Convolutional Autoencoder
  & From scratch
  & Pixel reconstruction error
  & In-domain encoder, multiple architectural variants \\
Variational Autoencoder (VAE)
  & From scratch
  & Reconstruction error + KL divergence
  & KL weight $\beta$ swept \\
Deep SVDD
  & From scratch
  & Distance to latent hypersphere centre
  & One-class geometric objective, no reconstruction \\
\midrule
Pretrained backbone (center-distance)
  & None (frozen)
  & Global embedding distance to normal centroid
  & ResNet18, WideResNet50-2 \\
\midrule
Teacher-student distillation
  & Student trained on normals
  & Spatial discrepancy map (patch-level)
  & ResNet18, ResNet50, WideResNet50-2 \\
PatchCore
  & None (memory bank only)
  & Patch nearest-neighbour distance
  & ResNet18, ResNet50, WideResNet50-2, EfficientNet-B0, ViT-B/16 \\
FastFlow
  & Normalising flow trained on normals
  & Log-likelihood under learned feature distribution
  & WideResNet50-2 (multilayer) \\
\bottomrule
\end{tabularx}
\end{table}

\subsection{Evaluation Protocol}

Every model uses the same deployment threshold derived from validation-normal
scores alone, with no access to defect labels:
\[
  \tau_{0.95} = \text{95th percentile of validation-normal anomaly scores.}
\]
\textbf{Deployed threshold.} The 95th percentile was chosen as a deliberate conservative operating point. By construction, exactly 5\% of validation normals exceed $\tau_{0.95}$, bounding the expected false-positive rate on normal wafers at approximately 5\% without requiring any defect labels for calibration. A stricter percentile (e.g.\ 99th) would reduce false alarms further but raise the threshold, causing more genuine defects to fall below it and be missed.

\textbf{Best-sweep F1.} This is a theoretical upper-bound metric that selects the threshold maximising F1 on the labelled test set. It quantifies the ceiling achievable through threshold tuning alone and is never used as the actual deployment threshold.

\begin{table}[H]
\centering
\caption{Metrics reported for each model. Val-threshold F1 is the primary metric.}
\label{tab:metrics}
\begin{tabularx}{\linewidth}{>{\raggedright\arraybackslash}p{3.2cm}
                              >{\raggedright\arraybackslash}X
                              >{\raggedright\arraybackslash}p{2cm}}
\toprule
Metric & Description & Role \\
\midrule
\textbf{Val-threshold F1} & Harmonic mean of precision and recall at $\tau_{0.95}$ & \textbf{Primary} \\
Precision @ $\tau_{0.95}$ & Fraction of flagged wafers that are truly defective & Supporting \\
Recall @ $\tau_{0.95}$ & Fraction of defective wafers successfully detected & Supporting \\
AUROC & Ranking quality across all thresholds & Ranking \\
AUPRC & Ranking quality under class imbalance (more informative than AUROC at 5\% defect rate) & Ranking \\
Best-sweep F1 & Best achievable F1 by sweeping thresholds on test labels & Ceiling only \\
\bottomrule
\end{tabularx}
\end{table}

Best-sweep F1 is never used as a deployment result. It quantifies how much
headroom remains if a more principled threshold rule were available. All models
are implemented in PyTorch (\texttt{torch}, \texttt{torchvision}, \texttt{timm},
\texttt{scikit-learn}, \texttt{umap-learn}). Full library versions are listed in \texttt{README.md}.

\subsection{Experimental Phases and Model Families}

The experiments proceed through four phases, each motivated by findings from the previous one.
Full training configurations and hyperparameter details are listed in Appendix~\ref{app:notebooks}.

\subsubsection{Phase 1: In-Domain Reconstruction Baselines}

\textbf{Motivation:} Establish a simple pipeline and identify what reconstruction-based
methods can and cannot do.

\textbf{Methods:} The convolutional autoencoder (latent dim 128, Adam $10^{-3}$, early stopping)
is trained from scratch on normal wafers only. A family of variants tests whether architectural
improvements can overcome the reconstruction paradigm's intrinsic limitations:
\begin{itemize}[leftmargin=1.5em]
  \item BatchNorm, dropout sweep $\{0.00, 0.05, 0.10, 0.20\}$, residual skip connections
  \item Resolution experiments: $64\!\times\!64$, $128\!\times\!128$, and $224\!\times\!224$
  \item Score ablation across seven reduction rules (MSE, MAE, max pixel, top-$k$ mean, foreground variants)
\end{itemize}
Two alternative one-class methods provide reference points: a Variational Autoencoder (VAE with $\beta$ swept)
and Deep SVDD (one-class hypersphere). A post-training score ablation reveals that scoring choice is
as impactful as architecture.

\textbf{Why this works:} Full-image reconstruction error is the simplest anomaly signal: normal wafers
reconstruct well, anomalous ones do not. This approach requires only normal training data and no pretrained models.

\textbf{Expected outcome:} Broad, globally visible defects such as Edge-Ring will be detectable through
elevated reconstruction error, but small, localised defects (Scratch, Loc) will be missed because they affect
only a few pixels in a 64$\times$64 or 224$\times$224 image. The fundamental limitation is that reconstruction error is a global signal.
Anomalies are sparse and local, but reconstruction-based scoring averages across the entire image, diluting
small defect signals. Overcoming this requires methods that preserve and score spatial information.

\subsubsection{Phase 2: Pretrained Features Without Local Scoring (Negative Control)}

\textbf{Motivation:} Test whether rich ImageNet-pretrained representations alone, without any spatial scoring, can
overcome the reconstruction ceiling.

\textbf{Methods:} Frozen ImageNet-pretrained ResNet18 and WideResNet50-2 backbones are evaluated using global
center-distance scoring: L2 distance from each test wafer's embedding to the centroid of train-normal embeddings.
No training or adaptation occurs. The backbones are used as-is.

\textbf{Why this is important:} This phase isolates whether the failure in Phase 1 is due to weak features or weak aggregation.
Pretrained backbones provide far richer features than in-domain autoencoders. If center-distance scoring succeeds with
these features, the bottleneck was feature quality. If it fails, the bottleneck is spatial averaging itself.

\textbf{Expected outcome:} Both baselines are expected to fall below Phase 1, confirming that feature quality alone
cannot overcome global-aggregation limitations and motivating the patch-level scoring methods in Phase 3.

\subsubsection{Phase 3: Local Spatial Scoring on Frozen Backbones}

\textbf{Motivation:} Having established that global aggregation fails, explore whether local patch-level anomaly
scoring on frozen pretrained features can outperform training-from-scratch.

\textbf{Methods:} Three methods test patch-level anomaly signals, each keeping the backbone frozen:

\begin{itemize}[leftmargin=1.5em]
  \item \textbf{Teacher-Student Distillation:} Train a lightweight student CNN to reproduce the frozen teacher's
    spatial feature maps on normal wafers. At inference, the per-location mismatch map is reduced to a wafer score
    by averaging the most deviant patch locations (top-k ratio swept, branch weights swept post-training).
    Backbones: ResNet18, ResNet50, WideResNet50-2 (single-layer and multilayer).

  \item \textbf{PatchCore:} Store a coreset-reduced memory bank of patch-level embeddings from a single forward
    pass on normal wafers. Score test wafers by computing nearest-neighbour distance of their patch embeddings
    to the bank, then aggregate per-patch anomaly scores to a wafer score. Training-free after feature extraction.
    Backbones: ResNet18, ResNet50, WideResNet50-2, EfficientNet-B0, EfficientNet-B1, Vision Transformer (ViT-B/16).

  \item \textbf{FastFlow:} Fit a normalising flow directly on top of frozen WideResNet50-2 spatial feature maps,
    estimating the feature distribution of normal wafers. Score test wafers by log-likelihood under the learned
    distribution. Requires training the flow head but not the backbone.
\end{itemize}

\textbf{Why local scoring matters:} Averaging small defects (50 pixels) with 4000+ normal pixels dilutes the signal;
better features cannot overcome this averaging bottleneck. Local patch-level scoring inverts the logic: each patch
is scored independently, preserving signals from small, localised defects while avoiding image-level averaging.

\textbf{Expected outcome:} If Phase 2 confirms that global aggregation is the bottleneck, then local scoring should
escape it. The three methods offer different trade-offs. Teacher-student enables post-training score tuning but
requires training a student head. PatchCore is training-free and memory-efficient with coreset reduction. FastFlow
requires training a flow head but avoids memory-bank storage. Each method tests whether local scoring alone can
overcome the global-aggregation limitation.


\subsubsection{Phase 4: Input Resolution and Backbone Scaling}

\textbf{Motivation:} If Phase 3 confirms that local scoring on pretrained features outperforms global methods, the next
question is whether resolution and backbone capacity can further improve performance. Phases 1--3 used either
upsampled or native $64\!\times\!64$ images. This phase tests the gain from preprocessing images at each backbone's
native resolution and whether stronger backbones provide additional benefits.

\textbf{Methods:} Wafer maps are natively only 32--64 pixels. Pretrained backbones expect $224\!\times\!224$ or higher.
In Phase 1--3, images were either upsampled inside the model or stored at $64\!\times\!64$. This phase tests whether
caching images at their backbone's native resolution (EfficientNet-B0: $224\!\times\!224$, EfficientNet-B1: $240\!\times\!240$,
WideResNet50-2: $224\!\times\!224$, ViT-B/16: $224\!\times\!224$) provides a lift, and whether scaling backbone capacity matters.
Each backbone is compared against its $64\!\times\!64$ counterpart with all other settings fixed.

For fair comparison, a plain autoencoder is also trained at native $224\!\times\!224$ resolution, allowing reconstruction-based
and pretrained-feature methods to be compared on equal footing.

Vision Transformer (ViT-B/16) is also qualitatively distinct: it processes the image as a sequence of non-overlapping
$16\!\times\!16$ patches with multi-head self-attention, producing globally context-aware patch features from the first layer.

\textbf{Why resolution and architecture matter:} When images are upsampled from 64 pixels to 224, spatial information
is interpolated, potentially blurring fine details. At native resolution, each pixel carries more information. Backbone
capacity (e.g., EfficientNet-B0 vs B1, or CNN vs Vision Transformer) determines what features can be learned or retrieved
from the pretrained ImageNet weights.

\textbf{Expected outcome:} Native resolution should improve performance by preserving pixel-level detail lost in interpolation.
Stronger backbones (e.g., EfficientNet-B1 vs B0, Vision Transformer vs CNN) should offer richer features. The key
question is whether these improvements are orthogonal to or limited by the local-scoring bottleneck identified in
Phases 1--2. Vision Transformers' self-attention mechanism may prove advantageous for capturing long-range spatial
relationships in small, localised defects.

\section{Main Results}

\subsection{Overview}

Table~\ref{tab:main_leaderboard} summarises ten representative models
that anchor the main narrative. Full sweep results across all variants are given in
Appendix~\ref{app:notebooks}.

\begin{table}[H]
\centering
\caption{Representative results on the main benchmark (ranked by val-threshold F1).
Models are selected to illustrate each stage of the experimental progression
rather than to list every variant tested.}
\label{tab:main_leaderboard}
\begin{tabularx}{\linewidth}{>{\raggedright\arraybackslash}Xcccc}
\toprule
Model & F1 & AUROC & AUPRC & Sweep F1 \\
\midrule
PatchCore + ViT-B/16, direct $224\!\times\!224$              & \textbf{0.595} & \textbf{0.956} & \textbf{0.671} & \textbf{0.651} \\
PatchCore + EfficientNet-B1, direct $240\!\times\!240$       & 0.591 & 0.935 & 0.609 & 0.651 \\
PatchCore + WideResNet50-2, direct $224\!\times\!224$        & 0.549 & 0.931 & 0.659 & 0.634 \\
PatchCore + WideResNet50-2, $64\!\times\!64$                 & 0.532 & 0.917 & 0.562 & 0.586 \\
Teacher-Student + ResNet50                                   & 0.525 & 0.909 & 0.599 & 0.606 \\
Teacher-Student + WideResNet50-2 (layer2+layer3)             & 0.524 & 0.923 & 0.546 & 0.561 \\
Autoencoder (plain, $224\!\times\!224$, native resolution)   & 0.510 & 0.901 & 0.596 & 0.587 \\
Autoencoder + BatchNorm (best score)                         & 0.502 & 0.834 & 0.568 & 0.630 \\
FastFlow + WideResNet50-2                                    & 0.482 & 0.871 & 0.489 & 0.482 \\
WideResNet50-2 center-distance (frozen)                      & 0.243 & 0.677 & 0.142 & 0.270 \\
ResNet18 center-distance (frozen)                            & 0.236 & 0.685 & 0.195 & 0.260 \\
\bottomrule
\end{tabularx}
\end{table}

\begin{figure}[H]
    \centering
    \includegraphics[width=\linewidth]{images/overall_experiment_comparison.png}
    \caption{Overall comparison across all completed experiments. Left: top runs
    ranked by deployed F1. Right: all runs in AUPRC-vs-F1 space, point size
    scaled by AUROC. The gap between the reconstruction family (lower left) and
    the local pretrained-feature methods (upper right) is the central finding
    of the project. Generated by \texttt{experiments/anomaly\_detection/report\_figures/notebook.ipynb}.}
    \label{fig:overall_comparison}
\end{figure}

Three cross-cutting findings run through every model family and explain the
shape of the leaderboard.

\textbf{1. Score design matters almost as much as model design.}
The BatchNorm autoencoder with its default scoring rule ranked below
the plain baseline ($\mathrm{F1}=0.431$ vs $0.468$). Post-training rescoring
with maximum absolute pixel error lifted it to $\mathrm{F1}=0.502$ without
any retraining. The same interaction appeared in the teacher-student family:
the ResNet50 teacher at its default score reached $\mathrm{F1}=0.488$. A
post-training score sweep lifted it to $\mathrm{F1}=0.525$. The scoring rule
is a design decision co-equal with the model architecture.

\textbf{2. Pretrained backbones require local scoring to be useful.}
Both center-distance baselines ($\mathrm{F1}<0.25$, $\mathrm{AUROC}<0.70$)
are weaker than every reconstruction model despite using far richer feature
representations. The bottleneck is not feature quality but global aggregation:
averaging all spatial information into one embedding vector discards exactly
the local patch evidence that defects produce.

\textbf{3. Native high-resolution preprocessing provides a clear and consistent lift.}
EfficientNet-B0 PatchCore rose from $\mathrm{F1}=0.467$ to $0.544$ when
wafer images were cached at full resolution rather than upsampled from
$64\!\times\!64$ inside the model. Scaling backbone capacity from B0 to B1 at
its native $240\!\times\!240$ input resolution pushes performance further to
$\mathrm{F1}=0.591$, making it the strongest CNN PatchCore result in the project.
WideResNet50-2 PatchCore improved similarly at $224\!\times\!224$, and
the ViT-B/16 branch achieves the overall best result ($\mathrm{F1}=0.595$)
at that same resolution.

\subsection{Phase 1 Results: Reconstruction Family}

\subsubsection{Reconstruction Family Performance Summary}

The autoencoder family achieved best result of $\mathrm{F1}=0.510$ at native $224\!\times\!224$ resolution, establishing the reconstruction-based ceiling. Across all variants, the pattern is consistent: broad defects are well-detected, small localised defects are missed. Key findings:

\begin{table}[H]
\centering
\small
\begin{tabularx}{\linewidth}{l X r r r r}
\toprule
\textbf{Variant} & \textbf{method} & \textbf{F1} & \textbf{AUROC} & \textbf{AUPRC} & \textbf{notes} \\
\midrule
AE x224 & topk\_abs\_mean & \textbf{0.510} & 0.901 & 0.596 & \textbf{best reconstruction} \\
AE+BN x64 & max\_abs & 0.502 & 0.834 & 0.568 & score ablation wins \\
AE baseline x64 & topk\_abs\_mean & 0.468 & 0.839 & 0.522 & baseline, strong scoring \\
AE+BN x64 & topk\_abs\_mean & 0.431 & 0.790 & 0.603 & same checkpoint, weak scoring \\
VAE x224 & $\beta=0.005$ & 0.340 & 0.772 & 0.362 & probabilistic baseline \\
Deep SVDD & hypersphere & 0.360 & 0.788 & 0.213 & one-class baseline \\
\bottomrule
\end{tabularx}
\caption{Phase 1 reconstruction family: all variants ranked by F1. Score choice is as impactful as architecture (BatchNorm at 0.431 vs 0.502 on same model). Resolution improvement (x64 → x224) is decisive.}
\label{tab:ae_family_summary}
\end{table}

\subsubsection{Architecture Design and Mechanism}

The baseline autoencoder follows a classic encoder-decoder design motivated by three core principles:

\begin{figure}[H]
\centering
\includegraphics[width=\linewidth]{images/autoencoder/architecture_diagram.png}
\caption{Convolutional autoencoder architecture: 64$\times$64 input $\rightarrow$ three strided convolution layers (16, 32, 64 channels) $\rightarrow$ flatten and fully connected layer to 128-dimensional latent space $\rightarrow$ symmetric decoder (fully connected, reshape, three transposed convolutions) $\rightarrow$ 64$\times$64 reconstruction via sigmoid. The information bottleneck at the latent layer forces the network to learn a compressed representation of normal wafers, which can then detect anomalies as elevated reconstruction error. Generated by \texttt{experiments/anomaly\_detection/autoencoder/x64/baseline/notebook.ipynb}.}
\label{fig:ae_architecture}
\end{figure}

\begin{enumerate}[leftmargin=1.5em]
  \item \textbf{Symmetric information bottleneck.} The encoder compresses 64$\times$64 spatial input through three strided convolutions (16, 32, 64 channels) into a 128-dimensional latent vector. This forces the network to discard all but the most essential patterns of normal wafers. The symmetry in the decoder (three transposed convolutions with matching strides) ensures that spatial relationships are inverted faithfully.

  \item \textbf{Progressive receptive field growth.} Each convolutional layer (stride=2, kernel 3$\times$3) doubles the receptive field while halving spatial dimensions. By the third layer, the receptive field spans most of the input, capturing global wafer structure. This hierarchical compression naturally learns multi-scale normal patterns (edges, defects, wafer geometry).

  \item \textbf{Information loss as anomaly detection.} The bottleneck forces reconstruction of normal wafers from just 128 values (99.6\% compression from 4096 pixels). Normal patterns, which the network has seen during training, compress and reconstruct well. Anomalies (defects and patterns absent from the training set) cannot be encoded compactly, so reconstruction error rises sharply. This is the signal we measure.
\end{enumerate}

The choice of 128-dimensional latent space balances two constraints: large enough to capture the diversity of normal wafers, small enough to force genuine compression. We do not tune this hyperparameter. 128 is a standard choice for image autoencoders at this scale.

\textbf{Alternative Architectures.} We test two reference methods that differ fundamentally from the reconstruction approach:

\begin{itemize}[leftmargin=1.5em]
  \item \textbf{Variational Autoencoder (VAE):} The encoder outputs a Gaussian mean and variance. A latent sample is drawn during training. The objective minimises reconstruction error plus a KL-divergence penalty ($\beta$ swept from 0.001 to 0.05). The best VAE result ($\mathrm{F1}=0.343$) underperforms the deterministic baseline ($\mathrm{F1}=0.468$), indicating that probabilistic regularisation does not benefit anomaly detection on this domain.

  \begin{table}[H]
  \centering
  \small
  \begin{tabularx}{\linewidth}{c X r r r}
  \toprule
  \textbf{KL weight} ($\beta$) & \textbf{role} & \textbf{F1} & \textbf{AUROC} & \textbf{AUPRC} \\
  \midrule
  0.001 & minimal regularisation & \textbf{0.343} & 0.770 & 0.336 \\
  0.005 & balanced & 0.340 & 0.772 & 0.372 \\
  0.010 & moderate regularisation & 0.335 & 0.766 & 0.369 \\
  0.050 & strong regularisation & 0.302 & 0.750 & 0.329 \\
  \bottomrule
  \end{tabularx}
  \caption{VAE KL-weight sweep: increasing $\beta$ strengthens the probabilistic prior but weakens reconstruction fidelity. All variants underperform the plain AE baseline (F1=0.468), confirming that the probabilistic objective does not help on this domain.}
  \label{tab:vae_beta_sweep}
  \end{table}

  \item \textbf{Deep SVDD:} Uses the same encoder architecture but trains a one-class geometric objective, minimising the volume of a latent hypersphere enclosing all normal embeddings. Wafers outside the hypersphere are flagged as anomalous. $\mathrm{F1}=0.360$, weaker than VAE, indicating that a single hypersphere cannot capture the diversity of normal wafers.
\end{itemize}

Both reference methods underperform plain reconstruction, which suggests that the reconstruction error signal, despite its global-averaging limitation, is fundamentally more informative than probabilistic or geometric alternatives on this dataset.

\subsubsection{Baseline Autoencoder Training and Score Selection}

\begin{figure}[H]
\centering
\includegraphics[width=0.95\linewidth]{images/autoencoder/ae64_training_and_evaluation.png}
\caption{Baseline AE: training curves (left) showing stable convergence by epoch 20, and score distribution (right) showing clear separation between normal and anomalous wafers. Generated by \texttt{experiments/anomaly\_detection/autoencoder/x64/baseline/notebook.ipynb}.}
\label{fig:ae64_training}
\end{figure}

The model converges within 20 epochs (best validation loss at epoch 38). The choice of reduction rule for aggregating per-pixel errors into a single wafer score substantially affects performance.

\begin{figure}[H]
\centering
\includegraphics[width=0.95\linewidth]{images/autoencoder/ae64_score_ablation.png}
\caption{Score ablation on the baseline AE checkpoint: seven reduction rules applied to the same checkpoint produce F1 ranging from 0.24 to 0.47. \texttt{topk\_abs\_mean} (top-20\% percentile of pixel errors) is strongest, suggesting that small clusters of high-error pixels carry more signal than full-image averages. Generated by \texttt{experiments/anomaly\_detection/autoencoder/x64/baseline/notebook.ipynb}.}
\label{fig:ae64_score_ablation}
\end{figure}

\textbf{Score Ablation:} Seven reduction rules were evaluated on the same fixed checkpoint:

\begin{table}[H]
\centering
\small
\begin{tabularx}{\linewidth}{l X r}
\toprule
\textbf{Scoring Rule} & \textbf{Definition} & \textbf{F1} \\
\midrule
\texttt{topk\_abs\_mean} & Average top 20\% of pixel errors (sorted by magnitude) & \textbf{0.468} \\
\texttt{mse\_mean} & Average squared pixel errors across entire image & 0.410 \\
\texttt{max\_abs} & Single worst (maximum absolute) pixel error & 0.324 \\
\texttt{foreground\_mse} & Average error in foreground-masked region & 0.317 \\
\texttt{mae\_mean} & Average absolute pixel errors (MAE) & 0.300 \\
\texttt{pooled\_mae\_mean} & MAE with spatial pooling & 0.293 \\
\texttt{foreground\_mae} & MAE in foreground region & 0.239 \\
\bottomrule
\end{tabularx}
\caption{Score ablation: seven reduction rules on the same AE checkpoint. \texttt{topk\_abs\_mean} wins by +0.057 F1 over \texttt{mse\_mean}, revealing that selecting high-error pixels outperforms averaging all pixels.}
\label{tab:score_ablation}
\end{table}

\texttt{topk\_abs\_mean} outperforms \texttt{mse\_mean} by $+0.057$ F1, confirming that anomaly signals concentrate in a small number of high-error pixels. Global averaging dilutes this. Top-$k$ selection preserves it.

\textbf{Baseline Result:} The plain autoencoder with \texttt{topk\_abs\_mean} scoring achieves
$\mathrm{F1}=0.468$, $\mathrm{AUROC}=0.839$, $\mathrm{AUPRC}=0.522$.
Per-defect analysis shows the characteristic pattern: broad, globally visible defects such
as Edge-Ring (recall 0.810) and Center (recall 0.720) are well-detected,
while small, localised defects such as Loc (recall 0.206) and Scratch
(recall 0.200) are largely missed. This defect-size dependence motivates every
subsequent modelling direction.

\subsubsection{Reconstruction Mechanism and Failure Analysis}

\textbf{How Reconstruction Error Becomes an Anomaly Signal:} The autoencoder compresses normal wafers efficiently.
Anomalous wafers (with unseen defect patterns) produce elevated reconstruction error.

\textbf{The Core Problem: Global Averaging.} Averaging pixel-level errors into a single score dilutes small defects.
If all 4096 pixels are averaged, a 50-pixel defect produces:
\[
\text{wafer score} = \frac{\text{50 high-error pixels} + \text{3946 low-error pixels}}{4096} \approx \text{weak signal}
\]
The defect is invisible to global scoring despite clear error signals (Figure \ref{fig:ae_reconstruction_failure}, magma heatmaps).
This reveals the core bottleneck: anomalies are local and sparse, but scoring is global and averaged.

\begin{figure}[H]
\centering
\begin{subfigure}[t]{\linewidth}
  \includegraphics[width=1\linewidth]{images/autoencoder/reconstruction_examples.png}
  \caption{Reconstruction mechanism (top=inputs, bottom=reconstructions): Normal wafers (left 3 pairs) reconstruct nearly perfectly, pixel-for-pixel matches. Anomalous wafers (right 3 pairs) show visible distortions (defect regions are blurred or smoothed away because the model has never seen these patterns during training).}
\end{subfigure}
\vspace{0.3cm}
\begin{subfigure}[t]{\linewidth}
  \includegraphics[width=\linewidth]{images/autoencoder/failure_examples_fn.png}
  \caption{False negatives in detail (6 columns: 2 samples side-by-side, each showing input | reconstruction | error map). Despite high reconstruction error visible in the magma heatmaps (rightmost columns), these small defects escape detection. The error signal exists, the bright magma colors indicate elevated pixel-level error, but global averaging dilutes it below the decision threshold.}
\end{subfigure}
\caption{Baseline autoencoder failure analysis: the reconstruction mechanism detects defects and produces error signals, but global averaging hides small anomalies. Generated by \texttt{experiments/anomaly\_detection/autoencoder/x64/baseline/notebook.ipynb}.}
\label{fig:ae_reconstruction_failure}
\end{figure}

\textbf{Root Cause:} Small defects produce elevated reconstruction error, but global averaging dilutes the signal.
A 50-pixel Scratch defect in a 4096-pixel image creates high error locally but vanishes when averaged.
Broad defects (Edge-Ring, Center) survive averaging because they affect hundreds of pixels. Defect size determines
detectability: small defects are missed, large ones are detected.

Per-defect recall across all eight defect types confirms this size gradient (Figure~\ref{fig:ae_per_defect_recall}).

\begin{figure}[H]
\centering
\includegraphics[width=\linewidth]{images/autoencoder/per_defect_recall_comparison.png}
\caption{Per-defect recall breakdown: broad defects (Edge-Ring 81\%, Center 72\%) are easily detected. Small localised defects (Scratch 20\%, Loc 21\%) are mostly missed. A Scratch defect affects only 50 pixels in a 4096-pixel image. When averaged with 3946 normal pixels, the signal drowns. Generated by \texttt{experiments/anomaly\_detection/autoencoder/x64/baseline/notebook.ipynb}.}
\label{fig:ae_per_defect_recall}
\end{figure}

\subsubsection{Architectural and Resolution Variants}

Two directions are evaluated to address the global-averaging limitation: architectural modifications and input resolution scaling.

\textbf{Architectural Modifications:} Standard techniques tested: BatchNormalization, dropout (0--0.20),
and residual skip connections. The question is whether these address the global-averaging bottleneck.

\textbf{Resolution Scaling:} We test a plain autoencoder at $224!\times!224$ resolution, pre-processed in a single nearest-neighbour resize from the raw wafer maps, matching the input size expected by ImageNet-pretrained backbones without a secondary bilinear upsample.

\textbf{Score Tuning First:} Rescoring the BatchNorm checkpoint with \texttt{max\_abs} (worst pixel error)
instead of \texttt{topk\_abs\_mean} improves F1 from 0.431 to 0.502 without retraining. Score choice is
as impactful as architecture.

\textbf{Resolution Impact:} Native $224\!\times\!224$ achieves F1=0.510 (vs 0.468 at 64×64). Spatial detail at
native resolution outweighs architectural modifications at lower resolution.

\begin{figure}[H]
\centering
\includegraphics[width=0.95\linewidth]{images/autoencoder/score_ablation_summary.png}
\caption{Score ablation on x224 checkpoint: \texttt{topk\_abs\_mean} remains dominant even at higher resolution, confirming that the pixel-selection insight (best pixels > average) holds across resolutions. Generated by \texttt{experiments/anomaly\_detection/autoencoder/x224/main/notebook.ipynb}.}
\label{fig:ae224_score_ablation}
\end{figure}

\begin{figure}[H]
\centering
\includegraphics[width=\linewidth]{images/autoencoder/autoencoder_family_comparison.png}
\caption{Autoencoder family ranking across all variants (ranked by F1). Key variants shown: AE baseline at 64$\times$64 (topk scoring), AE with BatchNorm at 64$\times$64 (max and topk scoring), Residual AE, and plain AE at 224$\times$224 native resolution. Resolution ($224\!\times\!224$ baseline at 0.510) decisively beats architectural modifications (BatchNorm, residual) at 64$\times$64 resolution. Score choice also matters: the same BatchNorm checkpoint ranges from F1=0.431 (topk) to F1=0.502 (max). Generated by \texttt{experiments/anomaly\_detection/report\_figures/notebook.ipynb}.}
\label{fig:ae_family}
\end{figure}

VAE ($\mathrm{F1}=0.340$) and Deep SVDD ($\mathrm{F1}=0.360$) both underperform the tuned reconstruction baseline
($\mathrm{F1}=0.468$) and the $224\!\times\!224$ variant ($\mathrm{F1}=0.510$), as visualised alongside
Phase 2 baselines in Figure~\ref{fig:compact_baselines}. Additional variants (dropout sweeps, residual encoder-decoder,
$128\!\times\!128$ resolution) are documented in the appendix. None exceeded the best reconstruction result.

\subsubsection{Phase 1 Summary: The Global-Averaging Bottleneck}

The reconstruction family conclusively demonstrates two insights:

\begin{enumerate}[leftmargin=1.5em]
  \item \textbf{Score design matters as much as architecture.} The same checkpoint ranges from F1=0.431 to F1=0.502 depending on the reduction rule. Selecting \texttt{max\_abs} (focus on worst pixels) over \texttt{topk\_abs\_mean} (average top 20\%) is a 16\% F1 improvement without retraining.

  \item \textbf{Resolution helps, but small defects still fail.} The best reconstruction result (AE x224, F1=0.510) achieves excellent recall on broad defects (Edge-Ring 90\%+) but remains weak on small defects (Scratch/Loc ~20\%). This is fundamentally because reconstruction error, no matter how well-scored, is a global signal: small anomalies drown when averaged across the entire image.
\end{enumerate}

This ceiling (F1=0.510) becomes the benchmark for Phase 2 and Phase 3: any improvement beyond reconstruction must fundamentally change how anomaly signals are aggregated (from global, full-image average, to local, patch-level scores).

\subsection{Phase 2 Results: Pretrained Backbone Baselines (Negative Control)}

Phase 2 evaluates whether richer pretrained representations alone can exceed the reconstruction ceiling,
using global center-distance scoring on frozen ImageNet-pretrained ResNet18 and WideResNet50-2 backbones.
Both backbones fall below the reconstruction baseline, confirming that global aggregation, not feature quality, is the bottleneck.

\subsubsection*{Center-Distance Scoring (Global Embedding Baseline)}

Each wafer's spatial feature map is spatially averaged to produce a global vector $\mathbf{f}_{\text{wafer}} \in \mathbb{R}^C$.
The anomaly score is the Euclidean distance to the train-normal centroid $\boldsymbol{\mu}$:
$\|\mathbf{f}_{\text{wafer}} - \boldsymbol{\mu}\|_2$. No training is required beyond computing $\boldsymbol{\mu}$ on normal data.

\begin{table}[H]
\centering
\small
\caption{Backbones tested in Phase 2 center-distance baseline. Both are ImageNet-pretrained and frozen. Only their feature quality differs.}
\label{tab:phase2_backbones}
\begin{tabularx}{\linewidth}{>{\raggedright\arraybackslash}p{2.0cm} >{\raggedright\arraybackslash}p{3.0cm} c >{\raggedright\arraybackslash}X}
\toprule
\textbf{Backbone} & \textbf{Depth} & \textbf{Feat. dim} & \textbf{Layer extraction} \\
\midrule
ResNet18 & 18 layers, 4 residual blocks & 512 & layer3 (stride 8, $8\!\times\!8$ resolution at 64px input) \\
WideResNet50-2 & 50 layers, wider bottlenecks (2$\times$), 4 residual blocks & 2048 & layer3 (stride 8, $8\!\times\!8$ resolution at 64px input) \\
\bottomrule
\end{tabularx}
\end{table}

\textbf{Shared setup.} Both backbones extract features from layer3 for comparable receptive fields. The key difference is capacity: ResNet18 produces 512-dimensional features versus 2048 for WideResNet50-2.

\textbf{Why capacity does not help here.} Spatial mean-pooling discards all local structure regardless of backbone depth. A localised defect contributes equally to the global vector as background noise, diluting the anomaly signal. Larger backbones do not fix this.

\begin{figure}[H]
\centering
\includegraphics[width=\linewidth]{images/compact_baseline_comparison.png}
\caption{Non-leading baselines: VAE, Deep SVDD, and plain embedding-distance
backbones. All fall well below the reconstruction-family ceiling, establishing
that both probabilistic one-class methods and raw pretrained features are
insufficient without local scoring. Generated by \texttt{experiments/anomaly\_detection/report\_figures/notebook.ipynb}.}
\label{fig:compact_baselines}
\end{figure}

Frozen ResNet18 ($\mathrm{F1}=0.236$, $\mathrm{AUROC}=0.685$) and
WideResNet50-2 ($\mathrm{F1}=0.243$, $\mathrm{AUROC}=0.677$) with
center-distance scoring both fall below the plain autoencoder
despite using much richer ImageNet features. The failure is not the
backbone but the aggregation: collapsing a full spatial feature map into
a single global embedding vector loses exactly the local deviation signals
that defects produce. Larger backbones do not fix this. Only local scoring
methods do. These results are therefore deliberate negative controls that
motivate the local scoring methods in the next two subsections.

\subsection{Phase 3 Results: Local Spatial Scoring on Frozen Backbones}

Phase 3 evaluates three patch-level anomaly scoring methods on frozen pretrained backbones: teacher-student
distillation, FastFlow, and PatchCore. Methods are ordered by decreasing training complexity.
Teacher-student requires training a student head. FastFlow trains a normalising flow. PatchCore requires no training.
Within each method, backbones are varied to isolate representation quality from scoring strategy.
Collectively, these results confirm that local patch-level scoring is the decisive factor for exceeding the reconstruction ceiling.

\subsubsection{Teacher-Student Distillation}

Teacher-student distillation addresses the local scoring problem by training
a lightweight student network to reproduce the teacher's spatial feature maps
on normal wafers only. At inference, the per-location discrepancy between
teacher and student forms a spatial anomaly map, which is then reduced to a
wafer score by averaging over the most deviant patch locations.

\begin{table}[H]
\centering
\small
\caption{Teacher-student distillation results across backbone variants. ViT-B/16 catastrophically fails, motivating the shift to PatchCore.}
\label{tab:ts_results}
\begin{tabularx}{\linewidth}{>{\raggedright\arraybackslash}X c c c c}
\toprule
\textbf{Backbone} & \textbf{F1} & \textbf{AUROC} & \textbf{AUPRC} & \textbf{Notes} \\
\midrule
ResNet18 (topk20) & 0.495 & 0.894 & 0.519 & Baseline \\
ResNet50 (mixed 2:1 topk20) & \textbf{0.525} & 0.909 & 0.599 & Best CNN result \\
WideResNet50-2 layer2 (topk25) & 0.508 & 0.920 & 0.540 & Single-layer \\
WideResNet50-2 multilayer (topk15) & 0.524 & \textbf{0.923} & 0.546 & Multilayer \\
ViT-B/16 (x224) & \textbf{0.163} & \textbf{0.661} & \textbf{0.106} & ⚠️ Architectural mismatch \\
\bottomrule
\end{tabularx}
\end{table}

ResNet18 with default scoring achieves $\mathrm{F1}=0.495$ after post-training score selection.
ResNet50 improves further. Default scoring gives $\mathrm{F1}=0.488$. A post-training branch-weight
sweep yields $\mathrm{F1}=0.525$, $\mathrm{AUROC}=0.909$, $\mathrm{AUPRC}=0.599$.
Scaling to WideResNet50-2 (single layer2) gives $\mathrm{F1}=0.508$, $\mathrm{AUROC}=0.920$.
Combining layer2 and layer3 gives $\mathrm{F1}=0.524$, $\mathrm{AUROC}=0.923$, $\mathrm{AUPRC}=0.546$,
matching ResNet50 on F1 with stronger ranking quality. Full sweep details are in the appendix.

\begin{figure}[H]
\centering
\includegraphics[width=\linewidth]{images/teacherstudent/ts_family_comparison_full.png}
\caption{Teacher-student distillation: all four backbone variants compared across
F1, AUROC, and AUPRC. WideResNet50-2 multilayer achieves the strongest AUROC
(0.923) among the teacher-student variants, confirming it as the best teacher
backbone. F1 is near-identical between ResNet50 and WRN50-2 multilayer (0.525
vs 0.524), but WRN50-2 provides clearly stronger anomaly ranking quality.
Generated by \texttt{experiments/anomaly\_detection/report\_figures/notebook.ipynb}.}
\label{fig:ts_family_full}
\end{figure}

\textbf{Note on ViT failure:} ViT-B/16 ($\mathrm{F1}=0.163$) is incompatible with teacher-student distillation.
ViT produces globally context-aware patch tokens via self-attention, which a CNN student cannot reproduce
through locally-growing receptive fields. This architectural mismatch motivated the shift to PatchCore,
which operates natively on patch embeddings without a learned student.

Teacher-student distillation also requires 25--30 epochs of student training and post-hoc branch-weight tuning,
raising the question of whether equivalent patch-level signals can be obtained without any training.

\subsubsection{FastFlow (Flow-Based Feature Density)}

FastFlow tests an alternative local density approach: rather than training a
student to mimic teacher features, it fits a normalising flow directly on top
of frozen WideResNet50-2 spatial feature maps, assigning anomaly scores through
log-likelihood under the learned normal feature distribution. This requires
training the flow head but not the backbone, and unlike PatchCore it does not
need a stored memory bank at inference.

\begin{figure}[H]
    \centering
    \includegraphics[width=\linewidth]{images/fastflow/fastflow_1.png}
    \caption{FastFlow preprocessing and feature extraction pipeline. Input grayscale wafer map is converted to 3-channel RGB, interpolated to backbone size, and passed through frozen WideResNet50-2. Features from layer2 (512 channels) and layer3 (2048 channels) are extracted and concatenated to form the combined feature tensor C, which feeds into the normalising flow. Generated by \texttt{experiments/anomaly\_detection/fastflow/x64/main/notebook.ipynb}.}
    \label{fig:fastflow_pipeline}
\end{figure}

\textbf{Stage 1: Feature extraction.} The input wafer map is converted from grayscale to RGB via channel repetition, resized to backbone input size via bilinear interpolation, then passed through a frozen WideResNet50-2 backbone. Features from layer2 (512 channels) and layer3 (2048 channels) are extracted and concatenated into a single tensor C (Figure~\ref{fig:fastflow_pipeline}).

\textbf{Stage 2: Normalising flow.} A flow with 4 sequential steps processes tensor C. Each step alternates between an orthogonal 1$\times$1 convolution (invertible channel mixing) and an affine coupling layer. The coupling layer splits channels in half, passes the first half through a small CNN (Conv3$\times$3 $\rightarrow$ GELU $\rightarrow$ Conv1$\times$1 $\rightarrow$ GELU $\rightarrow$ Conv3$\times$3, zero-initialised) to compute shift and log-scale parameters, then applies an affine transform to the second half.

\textbf{Stage 3: Anomaly scoring.} The transformed tensor Z is converted to a per-pixel anomaly map by computing $0.5 \times Z^2$ spatially. This heatmap is upsampled back to 64$\times$64 and reduced to a single wafer-level anomaly score via mean pooling.

\begin{figure}[H]
    \centering
    \includegraphics[width=\linewidth]{images/fastflow/fastflow_2.png}
    \caption{FastFlow normalising flow architecture. Each of 4 sequential steps processes feature tensor C through an orthogonal 1x1 convolution (invertible channel mixing) followed by an affine coupling layer. The affine coupling splits input channels, applies a learned transformation (parameterised by a small CNN shown in the diagram) to one half, and recombines. The final transformed tensor Z is converted to an anomaly map ($0.5 \times Z^2$) and reduced to a scalar wafer-level score. Inset: internal CNN of the affine coupling layer (Conv3x3 $\rightarrow$ GELU $\rightarrow$ Conv1x1 $\rightarrow$ GELU $\rightarrow$ Conv3x3, zero-initialised for identity initialisation). Generated by \texttt{experiments/anomaly\_detection/fastflow/x64/main/notebook.ipynb}.}
    \label{fig:fastflow_flow}
\end{figure}

\begin{table}[H]
\centering
\small
\caption{FastFlow configuration variants on WideResNet50-2 backbone.}
\label{tab:fastflow_results}
\begin{tabularx}{\linewidth}{>{\raggedright\arraybackslash}X c c c c}
\toprule
\textbf{Configuration} & \textbf{F1} & \textbf{AUROC} & \textbf{AUPRC} & \textbf{Notes} \\
\midrule
Layer2+Layer3, 4 steps, mean & \textbf{0.482} & \textbf{0.871} & \textbf{0.489} & Selected variant \\
Layer2+Layer3, 6 steps, mean & 0.470 & 0.870 & 0.479 & More flow depth unhelpful \\
Layer2 only, 6 steps, topk-r0.15 & 0.442 & 0.884 & 0.460 & Single-layer weaker \\
\bottomrule
\end{tabularx}
\end{table}

The selected multilayer configuration (layer2 + layer3, 4 flow steps, plain
mean reduction) achieves $\mathrm{F1}=0.482$, $\mathrm{AUROC}=0.871$,
$\mathrm{AUPRC}=0.489$. This is competitive with the reconstruction-family
ceiling ($\mathrm{F1}=0.502$) and demonstrates that flow-based density
modelling on pretrained features is a viable approach.

However, it falls clearly below the best teacher-student result ($\mathrm{F1}=0.525$) and
even further below the best WideResNet50-2 PatchCore result ($\mathrm{F1}=0.532$).
This performance gap suggests that while density modelling captures normal feature distributions,
it is less effective than explicit patch-level neighbourhood matching for detecting
the localised, feature-sparse anomalies in this dataset.

Per-defect recall reveals a distinctive pattern. FastFlow is relatively strong
on medium-scale localised defects, with Loc recall 0.588 and Edge-Loc recall 0.528.
In contrast, it is especially weak on the thin Scratch class (recall 0.133),
suggesting the learned density boundary does not capture the fine spatial geometry
of narrow streak defects.

These results position FastFlow as a meaningful reference point for the broader
anomaly detection landscape. However, they confirm that nearest-neighbour memory-bank
scoring (PatchCore) is a stronger approach on this dataset when the same backbone is used.

\begin{figure}[H]
\centering
\includegraphics[width=\linewidth]{images/fastflow/variant_comparison_metrics.png}
\caption{FastFlow variant comparison: F1, AUROC, AUPRC across three configurations. The multilayer $4$-step variant (layer2+layer3) outperforms both the 6-step multilayer and the single-layer configurations, showing that feature breadth matters more than flow depth. Generated by \texttt{experiments/anomaly\_detection/fastflow/x64/main/notebook.ipynb}.}
\label{fig:fastflow_variants}
\end{figure}

\begin{figure}[H]
\centering
\includegraphics[width=\linewidth]{images/fastflow/best_variant_defect_breakdown.png}
\caption{FastFlow per-defect recall for the selected variant (\texttt{wrn50\_l23\_s4}, mean reduction). Loc and Edge-Loc are the relative strengths, while Scratch (0.133) is the clear weak point, consistent with the broader project finding that thin linear defects are the hardest for all feature-density methods. Generated by \texttt{experiments/anomaly\_detection/fastflow/x64/main/notebook.ipynb}.}
\label{fig:fastflow_defect_breakdown}
\end{figure}

\subsubsection{PatchCore (Local Nearest-Neighbour Scoring)}

PatchCore stores a coreset-reduced memory bank of patch-level embeddings from normal wafers
and scores test wafers by nearest-neighbour distance to the bank. No training is required.
Patch-independent scoring directly targets localised defects and is natively compatible with ViT patch tokens.

\begin{table}[H]
\centering
\small
\caption{PatchCore results across backbone variants (ranked by F1).}
\label{tab:patchcore_results}
\begin{tabularx}{\linewidth}{>{\raggedright\arraybackslash}X c c c c}
\toprule
\textbf{Backbone} & \textbf{F1} & \textbf{AUROC} & \textbf{AUPRC} & \textbf{Notes} \\
\midrule
ViT-B/16 (224$\times$224) & \textbf{0.595} & \textbf{0.956} & \textbf{0.671} & Vision Transformer, best result \\
EfficientNet-B1 (240$\times$240) & 0.591 & 0.935 & 0.609 & Strongest CNN, native resolution \\
WideResNet50-2 (224$\times$224) & 0.549 & 0.931 & 0.659 & Multilayer features \\
EfficientNet-B0 (224$\times$224) & 0.544 & 0.925 & 0.483 & Native resolution lift \\
WideResNet50-2 (64$\times$64) & 0.532 & 0.917 & 0.562 & Upsampled input \\
ResNet50 (64$\times$64) & 0.420 & 0.821 & 0.363 & Medium backbone \\
ResNet18 (64$\times$64) & 0.401 & 0.842 & 0.411 & Shallow backbone \\
AE-BN encoder (64$\times$64) & 0.336 & 0.851 & 0.226 & In-domain features \\
\bottomrule
\end{tabularx}
\end{table}

The backbone variants are selected to isolate three factors: representation source (in-domain vs.\ ImageNet-pretrained), backbone capacity, and input resolution. Each step changes exactly one factor, allowing its contribution to be attributed directly.

\begin{figure}[H]
\centering
\includegraphics[width=\linewidth]{images/patchcore/patchcore_family_comparison.png}
\caption{PatchCore family: best result per backbone (left) and all PatchCore
variants coloured by backbone (right). The progression from in-domain AE
encoder to pretrained ResNet to WideResNet50-2 to ViT-B/16 follows a clear
upward trajectory, with native $224\!\times\!224$ preprocessing providing the
final decisive lift. Generated by \texttt{experiments/anomaly\_detection/report\_figures/notebook.ipynb}.}
\label{fig:patchcore_family}
\end{figure}

\begin{figure}[H]
\centering
\includegraphics[width=\linewidth]{images/patchcore/patchcore_backbone_progression.png}
\caption{PatchCore backbone progression ranked by val-threshold F1. The dashed
line marks the best reconstruction-family result ($\mathrm{F1}=0.502$).
Generated by \texttt{experiments/anomaly\_detection/report\_figures/notebook.ipynb}.}
\label{fig:patchcore_progression}
\end{figure}

Three design decisions account for most of the PatchCore progression.

\textbf{Backbone capacity} provides the largest gain ($+0.131$ F1): scaling from ResNet18 to WideResNet50-2 multilayer features at $64\!\times\!64$.

\textbf{Native resolution} contributes a further $+0.077$ F1 (EfficientNet-B0 x64 $\to$ x224), isolating the contribution of pixel-level detail independently of backbone changes. The jump from in-domain to ImageNet-pretrained features adds $+0.065$. Switching from CNN to ViT at the same native resolution provides a marginal additional $+0.004$.

PatchCore scales from in-domain features (F1=0.336) through ResNet18/50 to WideResNet50-2 multilayer (F1=0.532 at x64), establishing the Phase 3 ceiling. The remaining question is whether native-resolution preprocessing and stronger backbones can push beyond this point.

\subsection{Phase 4: Input Resolution and Backbone Scaling Results}

\textbf{Effect of input resolution.}
To isolate the contribution of input resolution independently of backbone capacity,
each backbone is evaluated at both $64\!\times\!64$ and its native resolution with all other settings fixed
(Figure~\ref{fig:resolution_comparison}).
EfficientNet-B0 provides the clearest controlled comparison: $\mathrm{F1}=0.467$ at $64\!\times\!64$
versus $\mathrm{F1}=0.544$ at native $224\!\times\!224$ ($+0.077$).
WideResNet50-2 PatchCore rises from $\mathrm{F1}=0.532$ to $0.549$ under the same resolution switch.
ViT-B/16 shows the steepest resolution dependence of any backbone tested
(Table~\ref{tab:vit_resolution}): $\mathrm{F1}$ collapses from $0.595$ at $224\!\times\!224$ to $0.342$
at $64\!\times\!64$ ($-0.253$).
The mechanism is architectural. At $64\!\times\!64$ the fixed $16\!\times\!16$ patch size produces only
$(64/16)^{2}=16$ tokens per image instead of $196$, so each token covers a region four times larger
in each spatial dimension. Defects that span fewer pixels than a single token become invisible to
the memory bank. CNN backbones degrade more gracefully because their convolutional receptive fields
scale proportionally with the input. ViT's token count is tied to the patch size, which is fixed.

\begin{table}[H]
\centering
\caption{Controlled resolution comparison across three PatchCore backbones. All other settings are
held fixed. Only the preprocessing resolution changes. ViT-B/16 suffers the largest absolute drop
because its token count falls from 196 to 16 at $64\!\times\!64$.}
\label{tab:vit_resolution}
\begin{tabular}{lcccc}
\toprule
Backbone & Resolution & F1 & AUROC & Tokens/image \\
\midrule
EfficientNet-B0 & $64\!\times\!64$   & 0.467 & 0.891 & N/A \\
EfficientNet-B0 & $224\!\times\!224$ & 0.544 & 0.925 & N/A \\
\midrule
WideResNet50-2  & $64\!\times\!64$   & 0.532 & 0.917 & N/A \\
WideResNet50-2  & $224\!\times\!224$ & 0.549 & 0.931 & N/A \\
\midrule
ViT-B/16        & $64\!\times\!64$   & 0.342 & 0.832 & 16 \\
ViT-B/16        & $224\!\times\!224$ & \textbf{0.595} & \textbf{0.956} & 196 \\
\bottomrule
\end{tabular}
\end{table}

\begin{figure}[H]
\centering
\includegraphics[width=0.75\linewidth]{images/resolution_effect_comparison.png}
\caption{Controlled resolution comparison for EfficientNet-B0 and WideResNet50-2
PatchCore. All other settings are held fixed. Only the wafer preprocessing
resolution changes. The gain is larger for EfficientNet-B0 ($+0.077$ F1) than
for WideResNet50-2 ($+0.017$ F1), suggesting EfficientNet-B0 is more sensitive
to the spatial detail captured at higher resolution. Generated by \texttt{experiments/anomaly\_detection/report\_figures/notebook.ipynb}.}
\label{fig:resolution_comparison}
\end{figure}

To understand the resolution improvement, we examine patch embedding geometry across our three strongest PatchCore variants using UMAP (Uniform Manifold Approximation and Projection). UMAP projects high-dimensional patch features into 2D while preserving local neighbourhood structure. Figure~\ref{fig:umap_3model} shows defect-type colouring for ResNet18 (F1=0.401), EfficientNet-B1 (F1=0.591), and ViT-B/16 (F1=0.595).

\textbf{Two distinct cluster patterns emerge.} EfficientNet-B1 produces the most compact within-defect-type clusters: defect points concentrate in tight blobs well separated from the normal mass. ViT-B/16, by contrast, scatters defect points across multiple peripheral satellite regions rather than a single cluster, yet achieves the highest F1 and AUROC on the main benchmark.

\textbf{Why ViT leads despite looser clusters.} The PatchCore anomaly score is the nearest-neighbour distance from each patch to the normal memory bank, not a measure of how tightly defects cluster among themselves. ViT's main-benchmark advantage arises from better overall normal/anomaly distance margins in high-dimensional space, which the 2D UMAP projection does not fully preserve. On the expanded holdout ($14\!\times$ more defects), EfficientNet-B1 overtakes ViT-B/16 (F1=0.596 vs 0.548, Section~\ref{sec:holdout}), suggesting that EfficientNet-B1's compact cluster geometry reflects more stable features that generalise better at scale.

\begin{figure}[H]
\centering
\includegraphics[width=\linewidth]{images/patchcore/umap_3model_defect_type.png}
\caption{UMAP patch embeddings coloured by defect type for three PatchCore backbones (grey = normal wafers). Generated by \texttt{experiments/anomaly\_detection/report\_figures/notebook.ipynb}.}
\label{fig:umap_3model}
\end{figure}

All three panels use a reference-fit UMAP trained on 5,000 train-normal embeddings only. Test points are projected via transform. ResNet18 and EfficientNet-B1 use Euclidean distance. ViT-B/16 uses cosine distance, so the three manifolds are not on the same metric space and should be read independently rather than compared geometrically.

\textbf{ResNet18 (left, F1=0.401):} Defect types are scattered throughout the normal cloud with no clear boundary.

\textbf{EfficientNet-B1 (centre, F1=0.591):} Defects form compact, well-separated blobs away from the normal mass.

\textbf{ViT-B/16 (right, F1=0.595):} Although defect points appear scattered across multiple peripheral regions in 2D, this backbone achieves the highest AUROC (0.956) of the three. The apparent inconsistency is resolved by recalling that PatchCore assigns each patch an anomaly score based on its distance to the nearest neighbour in the normal memory bank within the original high-dimensional feature space, not based on how compactly defects cluster in the 2D projection. ViT-B/16 achieves a stronger separation between defect and normal patches in that high-dimensional space than either ResNet18 or EfficientNet-B1, which is what drives its superior ranking performance.

\subsubsection{EfficientNet-B1 PatchCore: Strongest CNN Result}

Scaling from EfficientNet-B0 to EfficientNet-B1 at its native
$240\!\times\!240$ input resolution, with one-layer feature extraction from
block 3 and a coreset memory bank of 240k patches, pushes the CNN PatchCore
branch to its best result: $\mathrm{F1}=0.591$, $\mathrm{AUROC}=0.935$,
$\mathrm{AUPRC}=0.609$, Sweep F1 $=0.651$. This result is only 0.004 F1
below the ViT-B/16 leader and represents the strongest CNN-based
operating point in the project.

\begin{figure}[H]
\centering
\begin{subfigure}[t]{\linewidth}
  \includegraphics[width=\linewidth]{images/patchcore/effb1/benchmark_distribution_sweep.png}
  \caption{Score distribution and threshold sweep.}
\end{subfigure}
\vspace{0.3cm}
\begin{subfigure}[t]{0.38\linewidth}
  \includegraphics[width=\linewidth]{images/patchcore/effb1/benchmark_distribution_sweep_confusion.png}
  \caption{Confusion matrix at $\tau_{0.95}$: 195 TP, 55 FN, 215 FP, 4{,}785 TN.}
\end{subfigure}
\hfill
\begin{subfigure}[t]{0.58\linewidth}
  \includegraphics[width=\linewidth]{images/patchcore/effb1/benchmark_defect_breakdown.png}
  \caption{Per-defect recall breakdown.}
\end{subfigure}
\caption{EfficientNet-B1 PatchCore at $240\!\times\!240$: score distribution with deployed threshold (top left), confusion matrix at the 95th-percentile validation threshold (top right), and per-defect recall (bottom). Edge-Ring, Donut, Random, and Near-full all reach recall $\geq 0.93$, while Scratch (0.40) and Loc (0.59) remain the hardest categories. The confusion matrix shows 195 true positives out of 250 defects detected at a false-positive rate of 4.3\% (215 of 5{,}000 normals flagged).
Generated by \texttt{experiments/anomaly\_detection/patchcore/efficientnet\_b1/x240/main\_one\_layer/notebook.ipynb}.}
\label{fig:effnetb1_x240}
\end{figure}

\begin{table}[H]
\centering
\small
\caption{EfficientNet-B1 PatchCore per-defect recall at $240\!\times\!240$ vs.\ ViT-B/16 PatchCore. Defects ordered from hardest to easiest.}
\label{tab:effb1_per_defect}
\begin{tabular}{lccl}
\toprule
\textbf{Defect type} & \textbf{EffNet-B1} & \textbf{ViT-B/16} & \textbf{Gap} \\
\midrule
Scratch   & 0.400 & 0.727 & ViT $+0.327$ \\
Loc       & 0.588 & 0.829 & ViT $+0.241$ \\
Center    & 0.700 & 0.618 & \textbf{EffNet $+0.082$} \\
Edge-Loc  & 0.792 & 0.705 & \textbf{EffNet $+0.087$} \\
Edge-Ring & 0.929 & 0.941 & ViT $+0.012$ \\
Donut     & 1.000 & 1.000 & Tied \\
Random    & 1.000 & 1.000 & Tied \\
Near-full & 1.000 & 1.000 & Tied \\
\bottomrule
\end{tabular}
\end{table}

EfficientNet-B1 is the stronger CNN-based model and the expanded-holdout leader. Its per-defect profile reveals a clear structural pattern: it leads ViT-B/16 on broad, boundary-anchored defects (Center $+0.082$, Edge-Loc $+0.087$) but trails substantially on small localised defects (Scratch $-0.327$, Loc $-0.241$).

\textbf{Why EfficientNet-B1 leads on broad defects.} Its features form tight, geometrically stable clusters well-separated from the normal mass, as seen in the UMAP, which benefits defects with strong spatial contrast.

\textbf{Why ViT-B/16 leads on small defects.} Thin streaks and off-centre patches spanning only a few pixels require long-range spatial context. ViT's self-attention provides this from block 1, whereas convolutional receptive fields grow gradually and cannot resolve such fine details at the same scale.

\textbf{Holdout reversal.} On the expanded holdout, the per-class trade-off averages out in EfficientNet-B1's favour (F1=0.596 vs 0.548), suggesting that Scratch and Loc are rare enough that ViT's advantage on them does not outweigh EfficientNet-B1's broader geometric stability across a larger defect pool.

\subsubsection{ViT-B/16 PatchCore: Best Overall Result}

\textbf{Architecture.}
Vision Transformer (ViT-B/16) replaces convolutional feature extraction with a pure attention-based design.
The input is processed in three stages.

\textbf{Tokenization.}
A $224\!\times\!224$ wafer map is partitioned into $14\!\times\!14 = 196$ non-overlapping $16\!\times\!16$ patches.
Each patch is linearly projected to a $768$-dimensional token.
A learnable \texttt{[CLS]} token and additive positional embeddings are prepended, giving a sequence of $197$ tokens.

\textbf{Encoding.}
The sequence passes through $L=12$ identical encoder blocks.
Each block applies multi-head self-attention ($H=12$ heads, head dimension $64$) followed by a position-wise MLP ($768\!\to\!3072\!\to\!768$, GELU), both with residual connections.
Because self-attention operates jointly over all $197$ tokens at every block, each token has global receptive field from block~1, unlike CNNs where receptive fields grow gradually through strided convolutions.

\textbf{PatchCore integration.}
The $196$ patch-token embeddings from block~6 serve as per-patch feature vectors, stored in a greedy-coreset memory bank at a 5\% subsampling ratio (${\approx}400\text{k}$ patches) following \cite{roth2022patchcore}. The \texttt{[CLS]} token is discarded.
At inference, cosine distance from each patch embedding to its nearest memory-bank neighbour gives the per-patch anomaly score. The image-level score is the maximum over all $196$ patches.

Figure~\ref{fig:vit_architecture} shows the full pipeline.

\begin{figure}[H]
\centering
\includegraphics[width=\linewidth]{images/vit/vit_1.png}
\caption{ViT-B/16 PatchCore pipeline: patch tokenization, Transformer encoding, and coreset memory bank construction from layer-6 patch embeddings.}
\label{fig:vit_architecture}
\end{figure}

Applied with PatchCore at $224\!\times\!224$, ViT-B/16 achieves the best result in the project:
$\mathrm{F1}=0.595$, $\mathrm{AUROC}=0.956$, $\mathrm{AUPRC}=0.671$.
A controlled ablation at $64\!\times\!64$ ($\mathrm{F1}=0.342$, Table~\ref{tab:vit_resolution})
confirms that native $224\!\times\!224$ preprocessing is a hard requirement for this architecture.

\begin{figure}[H]
\centering
\begin{subfigure}[t]{\linewidth}
  \includegraphics[width=\linewidth]{images/vit/main/test_evaluation.png}
  \caption{Score distribution at the 95th-percentile validation threshold, threshold sweep
showing the headroom between the deployed and best-sweep operating points,
and confusion matrix at $\tau_{0.95}$.}
\end{subfigure}
\vspace{0.3cm}
\begin{subfigure}[t]{\linewidth}
  \includegraphics[width=\linewidth]{images/vit/main/defect_breakdown.png}
  \caption{Per-defect recall across all eight defect types.}
\end{subfigure}
\caption{ViT-B/16 PatchCore at direct $224\!\times\!224$, the best result in the
project ($\mathrm{F1}=0.595$, $\mathrm{AUROC}=0.956$, $\mathrm{AUPRC}=0.671$).
Generated by \texttt{experiments/anomaly\_detection/patchcore/vit\_b16/x224/main/notebook.ipynb}.}
\label{fig:vit_main_result}
\end{figure}

\begin{figure}[H]
\centering
\includegraphics[width=\linewidth]{images/vit/main/patchcore_heatmap_examples.png}
\caption{PatchCore (ViT-B/16, $224\!\times\!224$) per-patch anomaly localisation across all eight defect types.
Top row: original wafer binary maps (red pixels indicate failed dies, grey indicates empty dies).
Bottom row: per-patch anomaly score heatmap at the ViT patch grid resolution ($14\!\times\!14$), upsampled
to the native wafer resolution. Each heatmap is computed as the cosine distance from each patch embedding
to its nearest neighbour in a 500-normal mini memory bank. Bright regions indicate patches furthest from
the normal distribution. Defect regions are clearly highlighted even for small, spatially sparse defects
(Scratch, Loc) where image-level scoring fails. Generated by
\texttt{experiments/anomaly\_detection/patchcore/vit\_b16/x224/main/notebook.ipynb}.}
\label{fig:patchcore_heatmap}
\end{figure}

\textbf{Block depth sweep.}
A controlled ablation swept four extraction depths across the 12-block encoder
(blocks 3, 6, 9, and 11) with all other hyperparameters fixed.
Table~\ref{tab:vit_block_sweep} reports the results.

\begin{table}[H]
\centering
\caption{ViT-B/16 PatchCore block depth sweep at $224\!\times\!224$. All runs use
identical hyperparameters (5\% coreset, topk\_patch\_ratio~0.10, cosine distance).
Block indices are 0-based over the 12-block encoder.}
\label{tab:vit_block_sweep}
\begin{tabular}{llcccc}
\toprule
Block & Depth & F1 & AUROC & AUPRC & Recall \\
\midrule
3  & Early     & 0.550 & 0.934 & 0.622 & 0.732 \\
\textbf{6}  & \textbf{Mid}  & \textbf{0.580} & \textbf{0.954} & 0.642 & \textbf{0.800} \\
9  & Mid-late  & 0.580 & 0.943 & \textbf{0.677} & 0.788 \\
11 & Final     & 0.435 & 0.866 & 0.490 & 0.564 \\
\bottomrule
\end{tabular}
\end{table}

\textbf{Mid-depth blocks (6, 9) are optimal.} Both achieve $\mathrm{F1}=0.580$. Block~6 was selected for its higher AUROC and recall. Block~3 is competitive ($\mathrm{F1}=0.550$), showing that useful patch features accumulate early.

\textbf{The final block (11) collapses.} $\mathrm{F1}$ drops to $0.435$: by block~11, self-attention has pooled all patch representations into globally class-discriminative embeddings that lose the fine-grained spatial diversity PatchCore requires. This is consistent with findings that intermediate ViT layers retain the strongest local structure for dense prediction \cite{caron2021dino}.

\textbf{ViT gains are concentrated on small, localised defects.} At block~6, Scratch recall rises from $0.400$ to $0.727$ and Loc from $0.588$ to $0.829$ over EfficientNet-B1 (Table~\ref{tab:effb1_per_defect}). Because self-attention operates over all tokens from block~1, each patch embedding carries global spatial context, enabling detection of thin streaks and off-centre clusters that convolutional receptive fields cannot resolve at the same scale.

\textbf{ViT does not lead on every defect class.} It trails EfficientNet-B1 on Edge-Loc ($0.705$ vs $0.792$) and Center ($0.618$ vs $0.800$), both spatially broad defects where convolutional geometry provides sufficient signal without long-range attention.

These per-class trade-offs become clearer when placed in the context of the full experimental progression, from reconstruction baselines through to the best local-feature models.

\subsubsection{Cross-Model Defect Analysis}

Aggregate metrics such as F1 and AUROC summarise a model's overall discriminative power but conceal which defect morphologies it handles well or poorly.
A side-by-side comparison of per-defect recall across six representative models serves two purposes.
First, it reveals which defect classes remain hard for every family, indicating an inherent difficulty unrelated to architecture.
Second, it shows whether the progression from pixel-level reconstruction to local pretrained-feature scoring shifts the per-class recall profile in a systematic direction.

Table~\ref{tab:per_defect} and Figure~\ref{fig:per_defect_chart} compare
per-defect recall across six representative models spanning the full
progression, from the reconstruction baseline to the best local feature result.

\begin{table}[H]
\centering
\caption{Per-defect recall for six representative models tracing the experimental
progression. Bold indicates the best recall for that defect type.}
\label{tab:per_defect}
\begin{tabular}{lcccccc}
\toprule
Defect type
  & Baseline AE
  & AE-BN
  & WRN center
  & TS-Res50
  & WRN PC $224$px
  & ViT-B/16 PC \\
\midrule
Scratch   & 0.200 & 0.200 & 0.200 & 0.333 & 0.667 & \textbf{0.727} \\
Loc       & 0.206 & 0.176 & 0.176 & 0.559 & 0.735 & \textbf{0.829} \\
Edge-Loc  & 0.434 & 0.396 & 0.132 & 0.491 & 0.623 & \textbf{0.705} \\
Center    & 0.720 & 0.700 & 0.140 & 0.720 & \textbf{0.780} & 0.618 \\
Donut     & 0.571 & 0.429 & 0.286 & 0.714 & \textbf{1.000} & \textbf{1.000} \\
Edge-Ring & 0.810 & 0.833 & 0.464 & 0.929 & 0.798 & \textbf{0.941} \\
Random    & 0.600 & 0.600 & 0.600 & \textbf{1.000} & \textbf{1.000} & \textbf{1.000} \\
Near-full & 0.500 & \textbf{1.000} & 0.000 & \textbf{1.000} & \textbf{1.000} & \textbf{1.000} \\
\bottomrule
\end{tabular}
\end{table}

\textbf{Global aggregation fails uniformly.} The center-distance baseline is the weakest model on every defect class, confirming that collapsing spatial features into a single embedding discards the local deviation signals that defects produce.

\textbf{Reconstruction variants do not lead on any class.} AE-BN improves over the baseline only on Near-full ($0.500 \to 1.000$), a globally dominant pattern where improved feature normalisation is sufficient. Every other defect class is tied or regressed, confirming that architectural modifications within the reconstruction paradigm cannot escape the global-averaging bottleneck.

\textbf{Local scoring dominates on small, spatially confined defects.} WRN PatchCore and ViT-B/16 PatchCore lead across Scratch, Loc, Edge-Loc, and Donut. ViT-B/16 gains $+0.06$ on Scratch, $+0.09$ on Loc, and $+0.08$ on Edge-Loc over WRN PatchCore, driven by its globally context-aware patch embeddings from block~1.

\textbf{No single model leads on all classes.} WRN PatchCore achieves the highest Center recall ($0.780$), where spatially broad defects are well captured by dense local matching without requiring long-range attention. ViT-B/16 leads on Edge-Ring ($0.941$) with teacher-student close behind ($0.929$), but trails on Center ($0.618$ vs $0.780$).

\textbf{Scratch is the hardest class across all models.} Even the best model (ViT-B/16, recall $0.727$) leaves nearly $30\%$ of Scratch defects undetected. Its thin, low-coverage spatial signature produces only a handful of anomalous patches, making it the primary remaining challenge for any patch-level method.

\begin{figure}[H]
\centering
\includegraphics[width=\linewidth]{images/per_defect_model_comparison.png}
\caption{Per-defect recall across the six key models, shown as a grouped bar
chart. The shaded background highlights the three hardest defect classes
(Scratch, Loc, Edge-Loc). The progression from reconstruction-based methods
to local pretrained-feature methods is most visible in those three categories:
Scratch recall rises from 0.20 (Baseline AE) to 0.727 (ViT-B/16 PatchCore),
and Loc recall rises from 0.21 to 0.829. The WRN center-distance baseline
(grey) confirms that pretrained features alone without local scoring perform
worse than any trained model. Generated by \texttt{experiments/anomaly\_detection/report\_figures/notebook.ipynb}.}
\label{fig:per_defect_chart}
\end{figure}

Figure~\ref{fig:score_distributions} examines why aggregate F1 remains moderate even as per-defect recall improves.

\begin{figure}[H]
\centering
\includegraphics[width=\linewidth]{images/score_distribution_comparison.png}
\caption{Score distributions for normal (blue) and anomaly (red) wafers across
four key models. The dashed line marks the fixed 95th-percentile validation threshold (deployed F1);
the dotted green line marks the threshold that maximises F1 across a full score sweep (best-sweep F1).
Both values are shown in each panel title. Generated by \texttt{experiments/anomaly\_detection/report\_figures/notebook.ipynb}.}
\label{fig:score_distributions}
\end{figure}

Across all four models, anomaly scores shift rightward as model quality improves, but the distributions continue to overlap near the threshold.
This overlap is the primary explanation for the moderate deployed F1 values: a substantial fraction of low-anomaly-score defects fall below the 95th-percentile threshold and are missed, while some normal wafers exceed it.
The dotted green line in each panel shows where the decision boundary would sit if the threshold were chosen to maximise F1 on the test set.
The gap is largest for AE-BN (deployed $0.431$ vs.\ best-sweep $0.630$), where the score distribution is wide and the 95th-percentile falls well short of the optimal boundary.
For the two PatchCore models the gap is smaller (${\sim}0.06$), reflecting tighter, better-separated distributions, but it does not close entirely because the distributions still overlap near the boundary.

\textbf{What these two thresholds represent.}
The 95th-percentile threshold is a deployment-grade rule: it is calibrated from validation normals alone, requires no defect labels at any stage, and guarantees approximately a 5\% false-positive rate on normal wafers by construction.
It is the only threshold available in a real production line where defect labels do not exist before the model is deployed.
The best-sweep threshold is a theoretical upper bound: it scans every possible cut-point on the test set (which does carry labels) and selects the one that maximises F1.
It cannot be used in practice because finding it requires already knowing which wafers are defective, but it shows the maximum F1 gain achievable through better threshold selection alone, with all model weights held fixed.

\textbf{Precision-recall trade-off.} Moving from the deployed to the best-sweep threshold does not simply improve both precision and recall at once. In most cases the optimal threshold sits higher (stricter) than the 95th-percentile: it flags fewer wafers as anomalous, which eliminates many false positives and sharply raises precision, but at the cost of missing some low-score defects that the looser threshold would have caught.

Table~\ref{tab:threshold_comparison} illustrates this with the AE-BN model, which has the largest gap.

\begin{table}[H]
\centering
\caption{AE-BN threshold comparison: deployed 95th-percentile vs.\ best-sweep.
The best-sweep threshold is stricter (higher score required to flag a wafer),
so it flags far fewer wafers overall, raising precision dramatically at the
cost of reduced recall.}
\label{tab:threshold_comparison}
\begin{tabular}{lccccc}
\toprule
Threshold rule & Threshold value & Flagged wafers & Precision & Recall & F1 \\
\midrule
95th-pct (deployed) & 0.533 & ${\sim}330$ & 0.346 & 0.572 & 0.431 \\
Best-sweep (upper bound) & 0.954 & 165          & 0.789 & 0.524 & 0.630 \\
\bottomrule
\end{tabular}
\end{table}

The deployed threshold flags roughly $330$ wafers, of which only $0.346$ are genuine defects.
The best-sweep threshold raises the cut-point to $0.954$, halving the number flagged to $165$
and nearly tripling the precision to $0.789$, but recall falls from $0.572$ to $0.524$ because
the defects that score between $0.533$ and $0.954$ are now missed.
This pattern is diagnostic: the AE-BN score distribution is wide and poorly separated, so the
95th-pct of normals lands in a region densely populated by both normals and low-score defects.
Any threshold in that region is a poor operating point.
Better threshold selection cannot fix this. Only improving the score distribution itself can.
The PatchCore models show a smaller gap precisely because their anomaly distributions shift further
rightward, leaving less mixing at the decision boundary regardless of where the threshold is placed.

\section{Expanded Holdout Validation}
\label{sec:holdout}

To test whether the strongest checkpointed models remain stable on a much larger
evaluation pool, a disjoint secondary holdout was constructed:

\begin{table}[H]
\centering
\caption{Main benchmark vs.\ expanded holdout dataset sizes.}
\label{tab:holdout_sizes}
\begin{tabular}{lcc}
\toprule
Subset & Main benchmark & Expanded holdout \\
\midrule
Train normals & 40{,}000 & 40{,}000 \\
Val normals   &  5{,}000 &  5{,}000 \\
Test normals  &  5{,}000 & 70{,}000 \\
Test defects  &    250   &  3{,}500 \\
\bottomrule
\end{tabular}
\end{table}

The defect pool expands by roughly 14$\times$, making rare-class recall far more
meaningful than in the main benchmark. The same train and validation normals are
kept fixed, so each reevaluation reuses the previously selected checkpoint and
its validation-derived threshold policy.

\begin{table}[H]
\centering
\caption{Defect-type scaling from benchmark to expanded holdout.}
\label{tab:holdout_defects}
\begin{tabular}{lcc}
\toprule
Defect type & Benchmark & Holdout \\
\midrule
Edge-Ring & 84  & 1{,}302 \\
Edge-Loc  & 53  & 739     \\
Center    & 50  & 603     \\
Loc       & 34  & 492     \\
Scratch   & 15  & 169     \\
Random    &  5  & 108     \\
Donut     &  7  &  71     \\
Near-full &  2  &  16     \\
\bottomrule
\end{tabular}
\end{table}

\begin{table}[H]
\centering
\caption{Expanded-holdout results for one representative per model family plus the two
overall best checkpoints (70\,k normal / 3.5\,k defect test set).
The Autoencoder family is excluded (see text).}
\label{tab:holdout_leaderboard}
\begin{tabularx}{\linewidth}{>{\raggedright\arraybackslash}Xcccc}
\toprule
Model / variant & F1 & AUROC & AUPRC & Sweep F1 \\
\midrule
PatchCore + EfficientNet-B1, direct $240{\times}240$ & 0.596 & 0.953 & 0.656 & 0.671 \\
PatchCore + ViT-B/16, direct $224{\times}224$       & 0.548 & 0.941 & 0.615 & 0.606 \\
Teacher-Student + ResNet50                          & 0.516 & 0.926 & 0.622 & 0.585 \\
FastFlow + WideResNet50                             & 0.507 & 0.890 & 0.542 & 0.542 \\
\bottomrule
\end{tabularx}
\end{table}

\textbf{PatchCore + EfficientNet-B1} leads all reevaluated models on deployed F1 (0.596), AUROC (0.953),
and AUPRC (0.656), reversing the main-benchmark ranking where ViT-B/16 held a narrow 0.004 F1 edge.
This aligns with the UMAP observation that EfficientNet-B1 produces tighter, more compact defect
clusters (Figure~\ref{fig:umap_3model}).

\textbf{PatchCore + ViT-B/16} drops more sharply to F1=0.548, suggesting its main-benchmark advantage
may reflect threshold calibration sensitivity rather than more robust feature geometry.

\textbf{FastFlow + WideResNet50} holds a deployed F1 of 0.507, close to its main-benchmark performance,
confirming that its normalising-flow score scale transfers across normal populations.

\textbf{Autoencoder (excluded).} The Autoencoder family is excluded from this table because its
validation-derived threshold does not transfer to the holdout normal distribution: when evaluated on
the 70\,k holdout normals, 90\% exceed the 95th-percentile threshold calibrated on the 5\,k validation
normals, collapsing deployed F1 to 0.097 despite a best-sweep F1 of 0.652 (AUROC 0.845). This is a
concrete illustration of the threshold instability discussed in Section~\ref{sec:thresholding}.

\section{Diagnostic Analysis: Score Overlap and Threshold Limitations}

\textbf{The gap.} All models maintain a consistent gap between their deployed F1 and the best F1 achievable with optimal threshold selection. Score histograms reveal substantial overlap between normal and anomaly score distributions near the operating threshold, even for the best models. The key question is whether this gap arises from weak feature quality, poor geometric separability, or limitations of the 95th-percentile threshold rule itself.

\textbf{Diagnostic tools.} This section uses two approaches to investigate. First, UMAP visualises the patch embedding geometry to check whether normal and anomaly wafers are geometrically separable. Second, threshold sweeps scan every possible cut-point on the labelled test set to measure how much performance could be recovered with a perfectly chosen threshold.

\subsection{Embedding Geometry Analysis via UMAP}

UMAP projects patch embeddings into 2D to examine whether normal and anomalous wafers are geometrically separable. Given the small, sparse nature of wafer defects, perfect separation is not expected. The diagnostic question is what mixing pattern accounts for the gap between deployed and best-sweep F1. Figure~\ref{fig:umap_comparison} shows binary split-coloured views for the global baseline and WRN50-2 PatchCore. Figure~\ref{fig:nn_distance_vit} shows the nearest-neighbour distance distribution for ViT-B/16 PatchCore, quantifying score overlap directly in the original feature space. The backbone progression coloured by defect type is in Figure~\ref{fig:umap_3model}.

\begin{figure}[H]
\centering
\begin{subfigure}[t]{0.48\linewidth}
  \includegraphics[width=\linewidth]{experiments/anomaly_detection/backbone_embedding/wide_resnet50_2/x64/baseline/artifacts/umaps/wideresnet50A_embedding_baseline/evaluation/plots/embedding_umap.png}
  \caption{WRN50-2 center-distance baseline: anomalies (red) are distributed
  uniformly throughout the same islands as normals. No geometric threshold can
  reliably separate them.}
  \label{fig:umap_wrn_baseline}
\end{subfigure}
\hfill
\begin{subfigure}[t]{0.48\linewidth}
  \includegraphics[width=\linewidth]{experiments/anomaly_detection/patchcore/wideresnet50/x224/multilayer_umap_followup/artifacts/patchcore-wideresnet50-multilayer-umap/topk_mb50k_r005_x224/plots/umap_by_split.png}
  \caption{WRN50-2 PatchCore $224\!\times\!224$: anomalies begin to concentrate
  on manifold boundaries, thinner branches, and side islands rather than the
  dominant normal core, healthier geometry than the baseline, but still mixed.}
  \label{fig:umap_wrn_patchcore}
\end{subfigure}
\caption{Split-coloured UMAP comparison: global center-distance baseline (left)
vs. the WideResNet50-2 PatchCore $224\!\times\!224$ checkpoint (right).
Anomaly wafers are shown in red, normals in blue/grey. The progression from
completely uniform mixing to boundary-concentrated anomalies reflects the gain
from switching to local patch-level scoring.
Left panel generated by \texttt{experiments/anomaly\_detection/backbone\_embedding/wide\_resnet50\_2/x64/baseline/notebook.ipynb};
right panel by \texttt{experiments/anomaly\_detection/patchcore/wideresnet50/x224/multilayer\_umap\_followup/notebook.ipynb}.}
\label{fig:umap_comparison}
\end{figure}

\begin{figure}[H]
\centering
\includegraphics[width=\linewidth]{experiments/anomaly_detection/patchcore/vit_b16/x224/main/artifacts/patchcore_vit_b16_5pct/main_5pct/plots/nn_distance_distribution.png}
\caption{Nearest-neighbour distance distribution for ViT-B/16 PatchCore at $224\!\times\!224$.
The anomaly score is the maximum cosine distance from any of the 196 patch embeddings to its nearest
neighbour in the coreset memory bank. All three populations are shown on the same axes: validation
normals (dark, used only to derive the 95th-percentile threshold), test normals (blue), and test
anomalies (orange). The dashed red line marks the deployed threshold $\tau_{0.95}$; the dotted green
line marks the best-sweep threshold. The overlap between test-normal and test-anomaly distributions
near $\tau_{0.95}$ explains the gap between deployed F1 (0.595) and best-sweep F1 (0.651):
low-scoring anomalies fall below the threshold and are missed, while high-scoring normals exceed it.
Generated by \texttt{experiments/anomaly\_detection/patchcore/vit\_b16/x224/main/notebook.ipynb}.}
\label{fig:nn_distance_vit}
\end{figure}

\textbf{What the diagnostics reveal.} Overlapping geometry reflects the problem structure: defects are small,
sparse patches that barely perturb global embeddings. Global center-distance scoring fails
(Figure~\ref{fig:umap_comparison}, left) because no geometric threshold separates mixed clusters.
Patch-level methods succeed by scoring each patch independently against the memory bank, and the
resulting score distributions (Figure~\ref{fig:nn_distance_vit}) show that test-normal and test-anomaly
scores still overlap near the deployed threshold even for the best model. The gap between deployed F1
(0.595) and best-sweep F1 (0.651) reflects threshold placement within this overlap region, not weak
feature quality.

\section{Thresholding Limitations and Discussion}

The 95th-percentile validation-normal threshold is fair, reproducible, and
requires no defect labels.  However, three consistent observations from the
diagnostics show its limitations:

\begin{itemize}[leftmargin=1.5em]
  \item \textbf{Score overlap.} Even the best models show substantial overlap
        between normal and anomaly score distributions.  The histogram for
        the WRN PatchCore x224 selected variant (Figure~\ref{fig:wrn_x224_diag})
        shows that the normal and anomaly distributions overlap near the
        operating threshold.
  \item \textbf{Geometric overlap.} UMAP figures show that anomalies are not
        cleanly peeled off from the normal manifold. They occupy boundaries
        and tails, but do not form an isolated cluster.
  \item \textbf{Sweep headroom.} Best-sweep F1 is systematically $0.06$--$0.08$
        above deployed-threshold F1 across the strongest models (e.g.\ WRN
        PatchCore x224: deployed 0.549, sweep 0.634).  This gap represents
        the threshold-selection headroom that a more principled rule could capture.
\end{itemize}

Possible directions to improve thresholding beyond the 95th-percentile rule:
\begin{itemize}[leftmargin=1.5em]
  \item Adaptive thresholding calibrated on validation score structure.
  \item Local-density or manifold-aware decision boundaries informed by embedding geometry.
        The current UMAP-space KNN follow-ups are not yet sufficient:
        for example, the ViT-B/16 branch keeps a very strong original wafer
        score while its 2-D UMAP-KNN score collapses to near-random ranking.
  \item Hybrid supervision: if any defect examples can be used for threshold tuning
        without contaminating training, they could inform a more calibrated operating point.
\end{itemize}

Not every main-report winner has been reevaluated on the expanded holdout yet.
The holdout bundle is therefore a stability check on the reevaluated subset,
not a universal replacement benchmark.

\section{Comparison with Published Work}

Two published works provide the closest available reference points for
situating these results within the literature on normal-only wafer anomaly
detection. Direct numerical comparison is qualified carefully below, as
evaluation protocols differ across papers.

\paragraph{Jo and Lee (2021): one-class classification on WM-811K.}
The most directly comparable published study is \cite{jo2021oneclass},
which applies one-class classification to the identical WM-811K dataset
under a normal-only training regime. The paper trains and evaluates three
models: Deep SVDD \cite{ruff2018deepsvdd}, a standalone Adversarial
Autoencoder (AAE), and a hybrid AAE with DSVDD prior (AAE-DSVDD). Their
experimental design mirrors the present project in the key respects that
matter for comparison: training uses only labelled-normal wafers, and
evaluation is image-level binary detection across the same eight defect
pattern classes.

Their reported results show that Deep SVDD achieves precision and accuracy
below 0.5 across most classes, with AAE providing moderate improvement and
AAE-DSVDD achieving the best F1 among the three. This trajectory is
consistent with the present work: our Deep SVDD result (F1 = 0.360)
confirms the weakness of the pure hypersphere objective, and our
reconstruction family (best F1 = 0.502) already surpasses the AAE-class
performance reported in that work. The best result in the present project,
ViT-B/16 PatchCore at direct $224\!\times\!224$ (F1 = 0.595,
AUROC = 0.956), substantially exceeds all models reported in
\cite{jo2021oneclass} on image-level ranking quality, extending their
baseline study with the local pretrained-feature methods they do not
consider.

\paragraph{Roth et al.\ (2022): PatchCore on MVTec AD.}
The PatchCore method used as the primary local scoring approach in this
project is introduced by \cite{roth2022patchcore}. In that paper, PatchCore
is evaluated on the MVTec AD benchmark, achieving image-level
AUROC near 99\% across industrial texture and object categories under a
training-free, normal-only regime. While WM-811K is not evaluated in the
original paper, the method transfers directly to this domain: patch-level
nearest-neighbour scoring against a coreset-reduced memory bank of
normal-image features is architecture-agnostic. The present project can
therefore be understood as the first systematic application of PatchCore to
wafer-map anomaly detection under a strict deployment-style threshold
constraint, evaluating five backbone variants across two input resolutions.

The gap between MVTec AD performance (AUROC $\approx$99\%) and the present
WM-811K results (best AUROC = 0.956) is expected and informative. Wafer
maps present a structurally harder problem than most MVTec categories:
defects are sparse, occupying only a small fraction of die area. The normal
wafer texture is nearly uniform, reducing the contrast between normal and
defective patches, and the strict 95th-percentile threshold calibration on
validation normals alone removes the threshold-tuning headroom that
boosts reported AUROC-only benchmarks. The remaining 0.044 AUROC gap and
the deployed F1 of 0.595 are therefore attributable to the domain
difficulty and the deployment threshold constraint rather than to
implementation quality.

\section{Conclusion}

This project demonstrates a clear and reproducible progression in wafer-map
anomaly detection, implemented entirely in PyTorch:

\begin{enumerate}[leftmargin=1.5em]
  \item Plain reconstruction models establish a viable pipeline and show that
        local-defect sensitivity is the main limitation.
  \item Better in-domain autoencoder variants and score selection improve
        performance without new training. The BatchNorm autoencoder scored with
        maximum absolute pixel error is the strongest reconstruction baseline
        ($\mathrm{F1}=0.502$).
  \item Pretrained backbones with global center-distance scoring are consistently
        weak regardless of backbone size.
  \item Local anomaly scoring methods (teacher-student distillation and
        PatchCore) produce the strongest results. The score composition
        choice interacts strongly with the checkpoint. Post-training score
        sweeps are as important as architecture choices.
  \item Native high-resolution preprocessing provides a consistent and material
        lift: EfficientNet-B0 gained $+0.077$ F1 from $64\!\times\!64$ to
        $224\!\times\!224$. EfficientNet-B1 at $240\!\times\!240$ reaches
        $\mathrm{F1}=0.591$, the strongest CNN result in the project.
  \item The best completed result on the main benchmark, ViT-B/16 PatchCore at
        direct $224\!\times\!224$, achieves $\mathrm{F1}=0.595$,
        $\mathrm{AUROC}=0.956$, and $\mathrm{AUPRC}=0.671$. It leads
        EfficientNet-B1 by only 0.004 F1, with the ViT's advantage concentrated
        in Scratch, Loc, and Edge-Loc recall.
\end{enumerate}

Strong pretrained backbones combined with patch-level nearest-neighbour scoring substantially outperform
reconstruction-based and globally aggregated alternatives. The remaining bottleneck is score-distribution
overlap and threshold calibration, pointing toward more principled operating-point selection as the primary avenue for further improvement.

\appendix
\newpage
\section{Appendix: Group Contributions}
\label{app:contributions}

\begin{tabularx}{\linewidth}{>{\raggedright\arraybackslash}p{3.5cm}>{\raggedright\arraybackslash}X}
\toprule
Member & Primary contributions \\
\midrule
Henry Lee Jun (1004219)
  & Autoencoder family (baseline, BatchNorm, dropout sweep, residual, 128px variant),
    score ablation framework, Deep SVDD, VAE beta sweep, teacher-student distillation
    (ResNet18 and ResNet50), WideResNet50-2 teacher-student variants, score ensemble
    study, report writing and LaTeX preparation. \\
Chia Tang (1007200)
  & PatchCore backbone family (EfficientNet-B0/B1/B2, ResNet18, ResNet50, and WideResNet50-2 variants), 
    ViT-B16 PatchCore variants (one-layer/two-layer, with and without defect-tuning), threshold-tuning and score-analysis workflow, 
    UMAP embedding diagnostics, and experiment/report documentation for final write-up. \\
Genson Low (1005931)
  & \\
\bottomrule
\end{tabularx}

\newpage
\section{Appendix: How to Reproduce Each Figure}
\label{app:figures}

The rubric requires that each figure in the report can be recreated independently.
The table below maps every figure to the notebook or script that generates it.
All notebooks load pre-trained checkpoints by default (\texttt{RETRAIN = False})
and regenerate all plots from saved artifacts without requiring GPU access.
Setting \texttt{RETRAIN = True} reruns the full training and evaluation pipeline.

\begin{tabularx}{\linewidth}{>{\raggedright\arraybackslash}p{0.38\linewidth}>{\raggedright\arraybackslash}X}
\toprule
Figure & Source notebook \\
\midrule
Defect gallery (Fig.~\ref{fig:defect_gallery})
  & \texttt{artifacts/report\_plots/defect\_type\_gallery.png}, generated by the dataset notebook at \texttt{data/dataset/x64/benchmark\_50k\_5pct/notebook.ipynb} \\
Overall comparison (Fig.~\ref{fig:overall_comparison})
  & \texttt{artifacts/report\_plots/overall\_experiment\_comparison.png}, generated by running the report-plot aggregation cell in any top-level analysis notebook \\
AE training and score ablation (Fig.~\ref{fig:ae_training_ablation})
  & \texttt{experiments/anomaly\_detection/autoencoder/x64/baseline/notebook.ipynb} \\
AE x224 score ablation (Fig.~\ref{fig:ae224_score_ablation})
  & \texttt{experiments/anomaly\_detection/autoencoder/x224/main/artifacts/autoencoder\_x224/plots/score\_ablation\_summary.png}, generated by the x224 notebook score ablation cell \\
AE reconstruction examples (Fig.~\ref{fig:ae_reconstruction})
  & \texttt{experiments/anomaly\_detection/autoencoder/x64/baseline/notebook.ipynb} \\
AE family comparison (Fig.~\ref{fig:ae_family})
  & \texttt{artifacts/report\_plots/autoencoder\_family\_comparison.png}, regenerated by the autoencoder baseline notebook after all variant artifacts are present \\
Compact baselines (Fig.~\ref{fig:compact_baselines})
  & \texttt{artifacts/report\_plots/compact\_baseline\_comparison.png}, requires VAE, SVDD, and backbone-embedding artifacts \\
Teacher-student family (Fig.~\ref{fig:ts_family_full})
  & \texttt{experiments/anomaly\_detection/report\_figures/notebook.ipynb} \\
FastFlow results (Fig.~\ref{fig:fastflow_results})
  & \texttt{experiments/anomaly\_detection/fastflow/x64/main/notebook.ipynb} \\
PatchCore family comparison (Fig.~\ref{fig:patchcore_family})
  & \texttt{artifacts/report\_plots/patchcore\_family\_comparison.png}, requires all PatchCore variant artifacts \\
EfficientNet-B1 results (Fig.~\ref{fig:effnetb1_x240})
  & \texttt{experiments/anomaly\_detection/patchcore/efficientnet\_b1/x240/main\_one\_layer/notebook.ipynb} \\
WRN50-2 x224 diagnostics (Fig.~\ref{fig:wrn_x224_diag})
  & \texttt{experiments/anomaly\_detection/patchcore/wideresnet50/x224/multilayer/notebook.ipynb} \\
ViT-B/16 main result (Fig.~\ref{fig:vit_main_result})
  & \texttt{experiments/anomaly\_detection/patchcore/vit\_b16/x224/main/notebook.ipynb} \\
PatchCore heatmap examples (Fig.~\ref{fig:patchcore_heatmap})
  & \texttt{experiments/anomaly\_detection/patchcore/vit\_b16/x224/main/notebook.ipynb} \\
WRN baseline UMAP (Fig.~\ref{fig:umap_wrn_baseline})
  & \texttt{experiments/anomaly\_detection/backbone\_embedding/wide\_resnet50\_2/x64/baseline/notebook.ipynb} \\
ResNet18 PatchCore UMAP (Fig.~\ref{fig:umap_resnet18_patchcore})
  & \texttt{experiments/anomaly\_detection/patchcore/resnet18/x64/holdout70k\_3p5k\_umap\_followup/notebook.ipynb} \\
WRN50-2 PatchCore x224 UMAP (Fig.~\ref{fig:umap_wrn_patchcore})
  & \texttt{experiments/anomaly\_detection/patchcore/wideresnet50/x224/multilayer\_umap\_followup/notebook.ipynb} \\
ViT PatchCore NN distance distribution (Fig.~\ref{fig:nn_distance_vit})
  & \texttt{experiments/anomaly\_detection/patchcore/vit\_b16/x224/main/notebook.ipynb} \\
\midrule
\multicolumn{2}{l}{\textit{Cross-model comparison figures -- all generated by
\texttt{experiments/anomaly\_detection/report\_figures/notebook.ipynb}}} \\
Per-defect recall chart (Fig.~\ref{fig:per_defect_chart})
  & \texttt{report\_figures/notebook.ipynb}, Plot 1 \\
Resolution effect comparison (Fig.~\ref{fig:resolution_comparison})
  & \texttt{report\_figures/notebook.ipynb}, Plot 2 \\
PatchCore backbone progression (Fig.~\ref{fig:patchcore_progression})
  & \texttt{report\_figures/notebook.ipynb}, Plot 3 \\
Score distribution overlay (Fig.~\ref{fig:score_distributions})
  & \texttt{report\_figures/notebook.ipynb}, Plot 4 \\
TS family full comparison (Fig.~\ref{fig:ts_family_full})
  & \texttt{report\_figures/notebook.ipynb}, Plot 5 \\
3-model backbone UMAP by defect type (Fig.~\ref{fig:umap_3model})
  & \texttt{report\_figures/notebook.ipynb}, Plot 9 \\
\bottomrule
\end{tabularx}

\newpage
\section{Appendix: Experiment Inventory}
\label{app:notebooks}

This appendix documents the broader search space behind the main narrative.
Table~\ref{tab:all_experiments} lists every training run completed in this
project. Subsequent tables break down sweep results within each family.

\subsection{Complete Experiment List}

\begin{longtable}{>{\raggedright\arraybackslash}p{4.8cm}
                  >{\raggedright\arraybackslash}p{1.4cm}
                  c c c}
\caption{All experiments run in this project (main benchmark, best selected score per run).
Runs marked with $\dagger$ are intermediate steps superseded by a later variant in the same family.}
\label{tab:all_experiments} \\
\toprule
Experiment & Resolution & F1 & AUROC & AUPRC \\
\midrule
\endfirsthead
\multicolumn{5}{c}{\textit{(continued from previous page)}} \\
\toprule
Experiment & Resolution & F1 & AUROC & AUPRC \\
\midrule
\endhead
\midrule
\multicolumn{5}{r}{\textit{continued on next page}} \\
\endfoot
\bottomrule
\endlastfoot

\multicolumn{5}{l}{\textit{Reconstruction family}} \\
AE baseline (\texttt{topk\_abs\_mean}) & 64$\times$64 & 0.468 & 0.839 & 0.522 \\
AE baseline, longer run (43 ep)$^\dagger$ & 64$\times$64 & 0.460 & 0.835 & 0.525 \\
AE + BatchNorm (\texttt{max\_abs}) & 64$\times$64 & 0.502 & 0.834 & 0.568 \\
AE + BatchNorm + Dropout 0.00 & 64$\times$64 & 0.487 & 0.851 & 0.617 \\
AE + BatchNorm + Dropout 0.05 & 64$\times$64 & 0.483 & 0.835 & 0.552 \\
AE + BatchNorm + Dropout 0.10 & 64$\times$64 & 0.484 & 0.845 & 0.570 \\
AE + BatchNorm + Dropout 0.20 & 64$\times$64 & 0.470 & 0.841 & 0.575 \\
Residual AE (\texttt{max\_abs}) & 64$\times$64 & 0.474 & 0.843 & 0.589 \\
AE baseline (\texttt{topk\_abs\_mean}) & 128$\times$128 & 0.431 & 0.815 & 0.455 \\
VAE ($\beta=0.001$) & 64$\times$64 & 0.343 & 0.770 & 0.336 \\
VAE ($\beta=0.005$) & 64$\times$64 & 0.340 & 0.772 & 0.372 \\
VAE ($\beta=0.01$) & 64$\times$64 & 0.335 & 0.766 & 0.369 \\
VAE ($\beta=0.05$) & 64$\times$64 & 0.302 & 0.750 & 0.329 \\
Deep SVDD & 64$\times$64 & 0.360 & 0.788 & 0.213 \\
\midrule

\multicolumn{5}{l}{\textit{Pretrained backbone center-distance baselines}} \\
ResNet18 center-distance & 64$\times$64 & 0.236 & 0.685 & 0.195 \\
WideResNet50-2 center-distance & 64$\times$64 & 0.243 & 0.677 & 0.142 \\
\midrule

\multicolumn{5}{l}{\textit{Teacher-student distillation}} \\
TS + ResNet18 (student-only topk20) & 64$\times$64 & 0.495 & 0.894 & 0.519 \\
TS + ResNet50 (mixed 2:1 topk20) & 64$\times$64 & 0.525 & 0.909 & 0.599 \\
TS + WideResNet50-2 Layer2 (topk25) & 64$\times$64 & 0.508 & 0.920 & 0.540 \\
TS + WideResNet50-2 Multilayer (topk15) & 64$\times$64 & 0.524 & 0.923 & 0.546 \\
\midrule

\multicolumn{5}{l}{\textit{FastFlow}} \\
FastFlow WRN50-2 layer2+layer3, 4 steps & 64$\times$64 & 0.482 & 0.871 & 0.489 \\
FastFlow WRN50-2 layer2+layer3, 6 steps$^\dagger$ & 64$\times$64 & 0.470 & 0.870 & 0.479 \\
FastFlow WRN50-2 layer2 only, 6 steps$^\dagger$ & 64$\times$64 & 0.442 & 0.884 & 0.460 \\
\midrule

\multicolumn{5}{l}{\textit{PatchCore}} \\
PatchCore + AE-BN encoder (\texttt{mean}) & 64$\times$64 & 0.336 & 0.851 & 0.226 \\
PatchCore + ResNet18 (\texttt{mean}) & 64$\times$64 & 0.401 & 0.842 & 0.411 \\
PatchCore + ResNet50 (\texttt{mean}) & 64$\times$64 & 0.420 & 0.821 & 0.363 \\
PatchCore + WideResNet50-2 (\texttt{topk r=0.10}) & 64$\times$64 & 0.532 & 0.917 & 0.562 \\
PatchCore + WideResNet50-2 (\texttt{topk r=0.05}) & 64$\times$64 & 0.525 & 0.921 & 0.564 \\
PatchCore + WideResNet50-2 (\texttt{topk r=0.15}) & 64$\times$64 & 0.526 & 0.912 & 0.549 \\
PatchCore + EfficientNet-B0 (\texttt{topk r=0.02})$^\dagger$ & 64$\times$64 & 0.467 & 0.905 & 0.489 \\
PatchCore + WideResNet50-2 (\texttt{topk r=0.05}) & 224$\times$224 & 0.549 & 0.931 & 0.659 \\
PatchCore + EfficientNet-B0 (\texttt{topk r=0.02}) & 224$\times$224 & 0.544 & 0.925 & 0.483 \\
PatchCore + EfficientNet-B1 (\texttt{topk r=0.03}) & 240$\times$240 & 0.591 & 0.935 & 0.609 \\
PatchCore + ViT-B/16 block 6 (\texttt{topk r=0.10}) & 224$\times$224 & \textbf{0.595} & \textbf{0.956} & \textbf{0.671} \\
\midrule

\multicolumn{5}{l}{\textit{Score ensemble}} \\
WRN PatchCore + TS-Res50 (max fusion) & 64$\times$64 & 0.529 & 0.916 & 0.611 \\

\end{longtable}

\subsection{Autoencoder-Family Leaderboard (Main Benchmark)}

\begin{table}[H]
\centering
\caption{Selected autoencoder-family results (main benchmark, best score per checkpoint).}
\label{tab:appendix_ae}
\begin{tabular}{lccc}
\toprule
Variant & F1 & AUROC & AUPRC \\
\midrule
AE-224-topk               & 0.510 & 0.901 & 0.596 \\
AE-64-BN-max              & 0.502 & 0.834 & 0.568 \\
AE-64-BN-DO0.00-max       & 0.487 & 0.851 & 0.617 \\
AE-64-BN-DO0.10-max       & 0.484 & 0.845 & 0.570 \\
AE-64-BN-DO0.05-max       & 0.483 & 0.835 & 0.552 \\
AE-64-Res-max             & 0.474 & 0.843 & 0.589 \\
AE-64-topk (43 ep)        & 0.460 & 0.835 & 0.525 \\
AE-64-BN-topk             & 0.431 & 0.790 & 0.603 \\
AE-128-topk               & 0.431 & 0.815 & 0.455 \\
\bottomrule
\end{tabular}
\end{table}

\subsection{PatchCore Leaderboard (Main Benchmark)}

\begin{table}[H]
\centering
\caption{Selected PatchCore results (main benchmark, best variant per backbone).}
\label{tab:appendix_patchcore}
\begin{tabular}{lccc}
\toprule
Variant & F1 & AUROC & AUPRC \\
\midrule
ViT-B16-x224-topk10-mb400k        & 0.595 & 0.956 & 0.671 \\
EffNetB1-x240-topk3-mb240k        & 0.591 & 0.935 & 0.609 \\
WideRes50-x224-topk-mb600k-r005   & 0.549 & 0.931 & 0.659 \\
EffNetB0-x224-topk-mb240k-r002    & 0.544 & 0.925 & 0.483 \\
WideRes50-topk-mb50k-r010         & 0.532 & 0.917 & 0.562 \\
WideRes50-topk-mb50k-r015         & 0.526 & 0.912 & 0.549 \\
WideRes50-topk-mb50k-r005         & 0.525 & 0.921 & 0.564 \\
EffNetB0-x64-topk-mb240k-r002     & 0.467 & 0.905 & 0.489 \\
Res50-mean-mb50k                  & 0.420 & 0.821 & 0.363 \\
Res18-mean-mb50k                  & 0.401 & 0.842 & 0.411 \\
AE-BN-mean-mb50k                  & 0.336 & 0.851 & 0.226 \\
\bottomrule
\end{tabular}
\end{table}

\subsection{AE-Backed PatchCore Sweep}

\begin{table}[H]
\centering
\caption{PatchCore variants built on the frozen AE-BN encoder.}
\label{tab:appendix_aebn_patchcore}
\begin{tabularx}{\linewidth}{lcc>{\raggedright\arraybackslash}X}
\toprule
Variant & F1 & AUROC & Notes \\
\midrule
\texttt{mean\_mb50k}       & 0.336 & 0.851 & Best AE-BN PatchCore result \\
\texttt{topk\_mb50k\_r010} & 0.185 & 0.809 & Top-k reduction, 50\,k bank \\
\texttt{topk\_mb50k\_r005} & 0.123 & 0.777 & Narrower top-k region \\
\texttt{max\_mb50k}        & 0.056 & 0.679 & Max reduction too brittle \\
\bottomrule
\end{tabularx}
\end{table}

\subsection{ResNet18 Teacher-Student Focused Ablation}

\begin{table}[H]
\centering
\caption{ResNet18 teacher-student focused layer and top-k ablation.}
\label{tab:appendix_ts_resnet18}
\begin{tabular}{lcccc}
\toprule
Variant & Layer & Top-k ratio & F1 & AUROC \\
\midrule
\texttt{ts\_resnet18\_layer2\_topk25} & layer2 & 0.25 & 0.494 & 0.890 \\
\texttt{ts\_resnet18\_layer2\_topk20} & layer2 & 0.20 & 0.494 & 0.889 \\
\texttt{ts\_resnet18\_layer2\_topk15} & layer2 & 0.15 & 0.483 & 0.889 \\
\texttt{ts\_resnet18\_layer1\_topk15} & layer1 & 0.15 & 0.411 & 0.824 \\
\texttt{ts\_resnet18\_layer1\_topk20} & layer1 & 0.20 & 0.387 & 0.817 \\
\bottomrule
\end{tabular}
\end{table}

\subsection{ResNet50 Teacher-Student Score Sweep}

\begin{table}[H]
\centering
\caption{ResNet50 teacher-student local score-sweep summary (top rows).}
\label{tab:appendix_ts_resnet50}
\begin{tabular}{lcccc}
\toprule
Score variant & F1 & AUROC & AUPRC & Sweep F1 \\
\midrule
\texttt{s2\_a1\_topk\_mean\_r0.20} (selected) & 0.525 & 0.909 & 0.599 & 0.606 \\
\texttt{s1\_a0.5\_topk\_mean\_r0.20}           & 0.525 & 0.909 & 0.599 & 0.606 \\
\texttt{s1\_a2\_topk\_mean\_r0.20}             & 0.523 & 0.904 & 0.570 & 0.565 \\
\texttt{s2\_a1\_topk\_mean\_r0.10}             & 0.506 & 0.913 & 0.584 & 0.563 \\
\bottomrule
\end{tabular}
\end{table}

\subsection{WideResNet50-2 Teacher-Student Score Sweep}

\begin{table}[H]
\centering
\caption{WideResNet50-2 multilayer teacher-student score-sweep (top rows).}
\label{tab:appendix_ts_wrn}
\begin{tabular}{lcccc}
\toprule
Score variant & F1 & AUROC & AUPRC & Sweep F1 \\
\midrule
\texttt{s2\_a1\_topk\_mean\_r0.15} (selected) & 0.524 & 0.923 & 0.546 & 0.561 \\
\texttt{s1\_a1\_topk\_mean\_r0.25}             & 0.518 & 0.926 & 0.538 & 0.558 \\
\texttt{s2\_a1\_topk\_mean\_r0.20}             & 0.512 & 0.922 & 0.543 & 0.561 \\
\texttt{s2\_a1\_topk\_mean\_r0.10}             & 0.511 & 0.923 & 0.544 & 0.556 \\
\bottomrule
\end{tabular}
\end{table}

\subsection{FastFlow Ablation Summary}

\begin{table}[H]
\centering
\caption{FastFlow training-variant ablation (WideResNet50-2 backbone).}
\label{tab:appendix_fastflow}
\begin{tabular}{lccccc}
\toprule
Training variant & Layers & Steps & Best score & F1 & AUROC \\
\midrule
\texttt{wrn50\_l23\_s4} (selected) & layer2+layer3 & 4 & mean       & 0.482 & 0.871 \\
\texttt{wrn50\_l23\_s6}            & layer2+layer3 & 6 & mean       & 0.470 & 0.870 \\
\texttt{wrn50\_l2\_s6}             & layer2 only   & 6 & topk-r0.15 & 0.442 & 0.884 \\
\bottomrule
\end{tabular}
\end{table}

\subsection{Score Ensemble Study}

\begin{table}[H]
\centering
\caption{Post-hoc score ensemble: WRN PatchCore + TS-Res50.}
\label{tab:appendix_ensemble}
\begin{tabular}{lcccc}
\toprule
Fusion variant & F1 & AUROC & AUPRC & Sweep F1 \\
\midrule
PatchCore only         & 0.532 & 0.917 & 0.562 & 0.586 \\
max pair (selected)    & 0.529 & 0.916 & \textbf{0.611} & 0.620 \\
mean 70/30 PatchCore   & 0.528 & 0.918 & 0.596 & 0.616 \\
mean 50/50             & 0.526 & 0.917 & 0.606 & 0.615 \\
TS only                & 0.525 & 0.909 & 0.599 & 0.606 \\
\bottomrule
\end{tabular}
\end{table}

The ensemble improved AUPRC notably (0.611 vs 0.562) but did not beat the
standalone PatchCore on deployed F1.

\subsection{Experiment-to-File Reference}

\begin{table}[H]
\centering
\caption{Key experiment-to-notebook mapping for reproducibility.
All paths are relative to the project root.}
\label{tab:appendix_notebooks}
\begin{tabularx}{\linewidth}{>{\raggedright\arraybackslash}X>{\raggedright\arraybackslash}X}
\toprule
Experiment & Notebook path \\
\midrule
Baseline AE (224$\times$224)
  & \texttt{experiments/anomaly\_detection/autoencoder/x224/main/notebook.ipynb} \\
Baseline AE
  & \texttt{experiments/anomaly\_detection/autoencoder/x64/baseline/notebook.ipynb} \\
BatchNorm AE
  & \texttt{experiments/anomaly\_detection/autoencoder/x64/batchnorm/notebook.ipynb} \\
BatchNorm + Dropout sweep
  & \texttt{experiments/anomaly\_detection/autoencoder/x64/batchnorm\_dropout/notebook.ipynb} \\
Residual AE
  & \texttt{experiments/anomaly\_detection/autoencoder/x64/residual/notebook.ipynb} \\
AE 128$\times$128
  & \texttt{experiments/anomaly\_detection/autoencoder/x128/baseline/notebook.ipynb} \\
VAE baseline
  & \texttt{experiments/anomaly\_detection/vae/x64/baseline/notebook.ipynb} \\
VAE beta sweep
  & \texttt{experiments/anomaly\_detection/vae/x64/beta\_sweep/notebook.ipynb} \\
Deep SVDD
  & \texttt{experiments/anomaly\_detection/svdd/x64/baseline/notebook.ipynb} \\
ResNet18 backbone baseline
  & \texttt{experiments/anomaly\_detection/backbone\_embedding/resnet18/x64/baseline/notebook.ipynb} \\
WRN50-2 embedding baseline
  & \texttt{experiments/anomaly\_detection/backbone\_embedding/wide\_resnet50\_2/x64/baseline/notebook.ipynb} \\
PatchCore + AE-BN backbone
  & \texttt{experiments/anomaly\_detection/patchcore/ae\_bn/x64/main/notebook.ipynb} \\
PatchCore + ResNet18
  & \texttt{experiments/anomaly\_detection/patchcore/resnet18/x64/main/notebook.ipynb} \\
PatchCore + ResNet50
  & \texttt{experiments/anomaly\_detection/patchcore/resnet50/x64/main/notebook.ipynb} \\
PatchCore + WRN50-2 x64
  & \texttt{experiments/anomaly\_detection/patchcore/wideresnet50/x64/main/notebook.ipynb} \\
PatchCore + WRN50-2 x224
  & \texttt{experiments/anomaly\_detection/patchcore/wideresnet50/x224/multilayer/notebook.ipynb} \\
PatchCore + WRN50-2 x224 UMAP
  & \texttt{experiments/anomaly\_detection/patchcore/wideresnet50/x224/multilayer\_umap\_followup/notebook.ipynb} \\
PatchCore + EfficientNet-B0 x224
  & \texttt{experiments/anomaly\_detection/patchcore/efficientnet\_b0/x224/main/notebook.ipynb} \\
PatchCore + EfficientNet-B1 x240
  & \texttt{experiments/anomaly\_detection/patchcore/efficientnet\_b1/x240/main\_one\_layer/notebook.ipynb} \\
PatchCore + ViT-B/16 x224
  & \texttt{experiments/anomaly\_detection/patchcore/vit\_b16/x224/main/notebook.ipynb} \\
TS + ResNet18
  & \texttt{experiments/anomaly\_detection/teacher\_student/resnet18/x64/main/notebook.ipynb} \\
TS + ResNet50
  & \texttt{experiments/anomaly\_detection/teacher\_student/resnet50/x64/main/notebook.ipynb} \\
TS + WRN50-2 layer2
  & \texttt{experiments/anomaly\_detection/teacher\_student/wideresnet50\_2/x64/layer2\_self\_contained/notebook.ipynb} \\
TS + WRN50-2 multilayer
  & \texttt{experiments/anomaly\_detection/teacher\_student/wideresnet50\_2/x64/multilayer\_self\_contained/notebook.ipynb} \\
FastFlow sweep
  & \texttt{experiments/anomaly\_detection/fastflow/x64/main/notebook.ipynb} \\
Score ensemble
  & \texttt{experiments/anomaly\_detection/ensemble/x64/score\_ensemble/notebook.ipynb} \\
Report comparison figures
  & \texttt{experiments/anomaly\_detection/report\_figures/notebook.ipynb} \\
\bottomrule
\end{tabularx}
\end{table}

\newpage
\begin{thebibliography}{9}

\bibitem{khatun2026cbam}
M.~R. Khatun, F.~A. Farid, S.~Dhar, M.~S. Islam, J.~Uddin, and H.~A. Karim,
\newblock ``CBAM-enhanced lightweight CNN for wafer map defect classification,''
\newblock \textit{Frontiers in Electronics}, vol.~7, 2026.
\newblock \url{https://doi.org/10.3389/felec.2026.1750707}

\bibitem{taha2025survey}
K.~Taha,
\newblock ``Observational and experimental insights into machine learning-based defect classification in wafers,''
\newblock \textit{Journal of Intelligent Manufacturing}, 2025.
\newblock \url{https://doi.org/10.1007/s10845-024-02521-0}

\bibitem{jo2021oneclass}
H.~Y. Jo and S.-W. Lee,
\newblock ``One-class classification for wafer map using adversarial autoencoder with DSVDD prior,''
\newblock \textit{arXiv preprint arXiv:2107.08823}, 2021.
\newblock \url{https://arxiv.org/abs/2107.08823}

\bibitem{roth2022patchcore}
K.~Roth, L.~Pemula, J.~Zepeda, B.~Sch\"{o}lkopf, T.~Brox, and P.~Gehler,
\newblock ``Towards total recall in industrial anomaly detection,''
\newblock in \textit{Proceedings of the IEEE/CVF Conference on Computer Vision and Pattern Recognition (CVPR)},
\newblock pp.\ 14318--14328, 2022.
\newblock \url{https://arxiv.org/abs/2106.08265}

\bibitem{dosovitskiy2021image}
A.~Dosovitskiy, L.~Beyer, A.~Kolesnikov, D.~Weissenborn, X.~Zhai, T.~Unterthiner,
M.~Dehghani, M.~Minderer, G.~Heigold, S.~Gelly, J.~Uszkoreit, and N.~Houlsby,
\newblock ``An image is worth 16x16 words: Transformers for image recognition at scale,''
\newblock in \textit{International Conference on Learning Representations (ICLR)}, 2021.
\newblock \url{https://arxiv.org/abs/2010.11929}

\bibitem{caron2021dino}
M.~Caron, H.~Touvron, I.~Misra, H.~J\'{e}gou, J.~Mairal, P.~Bojanowski, and A.~Joulin,
\newblock ``Emerging properties in self-supervised vision transformers,''
\newblock in \textit{Proceedings of the IEEE/CVF International Conference on Computer Vision (ICCV)},
\newblock pp.\ 9650--9660, 2021.
\newblock \url{https://arxiv.org/abs/2104.14294}

\bibitem{ruff2018deepsvdd}
L.~Ruff, R.~Vandermeulen, N.~Goernitz, L.~Deecke, S.~A. Siddiqui, A.~Binder, E.~M\"{u}ller, and M.~Kloft,
\newblock ``Deep one-class classification,''
\newblock in \textit{Proceedings of the 35th International Conference on Machine Learning (ICML)},
\newblock vol.\ 80, pp.\ 4393--4402, 2018.
\newblock \url{http://proceedings.mlr.press/v80/ruff18a.html}

\end{thebibliography}

\end{document}
